\documentclass[11pt]{article}

\usepackage[margin=1.15in]{geometry}
\usepackage{amsmath,amssymb,amsthm}
\usepackage{mathtools}
\usepackage{xcolor}
\usepackage{framed}
\usepackage{array}
\usepackage{booktabs}
\usepackage[colorlinks=true,linkcolor=blue!50!black,citecolor=blue!50!black,urlcolor=blue!50!black]{hyperref}

% ---------------------------------------------------------------
% Theorem environments
% ---------------------------------------------------------------
\newtheorem{theorem}{Theorem}[section]
\newtheorem{lemma}[theorem]{Lemma}
\newtheorem{proposition}[theorem]{Proposition}
\newtheorem{corollary}[theorem]{Corollary}
\newtheorem{claim}[theorem]{Claim}

\theoremstyle{definition}
\newtheorem{definition}[theorem]{Definition}
\newtheorem{remark}[theorem]{Remark}

% ---------------------------------------------------------------
% "Gap / review point" environment: numbered, framed, colored
% ---------------------------------------------------------------
\definecolor{gapred}{RGB}{160,20,20}
\definecolor{gapbg}{RGB}{253,242,242}
\newcounter{gapctr}
\newenvironment{gap}[1]{%
  \refstepcounter{gapctr}%
  \def\FrameCommand{\fboxsep=8pt\fcolorbox{gapred}{gapbg}}%
  \begin{framed}\noindent\textbf{\color{gapred}Gap / review point \thegapctr\ (#1).}\ }%
  {\end{framed}}

\newcommand{\LB}{\Lambda}
\newcommand{\eps}{\varepsilon}
\newcommand{\R}{\mathbb{R}}
\newcommand{\C}{\mathbb{C}}
\newcommand{\Z}{\mathbb{Z}}
\newcommand{\calP}{\mathcal{P}}
\newcommand{\calS}{\mathcal{S}}
\newcommand{\CapT}{\operatorname{Cap}}
\DeclareMathOperator{\Rea}{Re}
\DeclareMathOperator{\Ima}{Im}
\DeclareMathOperator{\sgn}{sgn}

\title{An exposition of a computer-assisted proof of the bound
$\Lambda\le 0.1787854$ for the de Bruijn--Newman constant}
\author{Authors TBD\thanks{Referee's exposition of an unpublished
proof by J.~Gomila (private repository
\texttt{dbn-lambda-01787854-candidate-audit}, ``sealed referee snapshot,''
2026). This document describes the argument for the purpose of
independent review. At the present stage of the review every
analytic lemma carries a complete verified proof and every
computer-assisted component has been independently reproduced (see
Section~\ref{sec:gaps}); a final overall certification of
correctness awaits the completion of the review.}}
\date{July 27, 2026}

\begin{document}
\maketitle

\begin{abstract}
We give a self-contained mathematical exposition of a claimed computer-assisted proof, due to J.~Gomila, of the bound
\[
\LB\le\frac{893927}{5000000}=0.1787854
\]
for the de Bruijn--Newman constant $\LB$, improving the bound
$\LB\le0.22$ of the Polymath~15 project and the current record
$\LB\le0.2$ of Platt--Trudgian. The argument is intended to be
unconditional: it instantiates the criterion of Polymath~15
[Theorem~1.2, arXiv:1904.12438v2] at the exact parameters
$X=6000000185827$, $t_0=129/800$, $y_0^2=87677/2500000$, supplying its
three hypotheses by, respectively, the Platt--Trudgian finite
verification of the Riemann hypothesis, a hybrid finite-scan/analytic-tail
nonvanishing theorem for $H_{t_0}$, and an interval-arithmetic barrier
certificate. The proof has not been peer reviewed. In the course of
the review recorded here, every analytic lemma has been given a
complete verified proof; all statements quoted from the published
literature have been checked against the arXiv sources of the cited
papers; and each of the seven computer-assisted components has been
independently reproduced, by programs reviewed line by line against
the corresponding statements, running on an independently built
toolchain and reading no stored data. The verification status is
summarized in Section~\ref{sec:gaps}; a final overall assessment is
deferred to the conclusion of the review.
\end{abstract}

\tableofcontents

% ===============================================================
\section{Introduction}
% ===============================================================

\subsection{The de Bruijn--Newman constant}

Let $\Phi$ denote the super-exponentially decaying function
\[
\Phi(u)=\sum_{n=1}^{\infty}
\bigl(2\pi^2n^4e^{9u}-3\pi n^2e^{5u}\bigr)
\exp\!\bigl(-\pi n^2e^{4u}\bigr),
\]
and for $t\in\R$ define the entire functions
\begin{equation}\label{eq:Ht}
H_t(z)=\int_0^\infty e^{tu^2}\Phi(u)\cos(zu)\,du .
\end{equation}
The function $H_0$ is a normalization of the Riemann $\xi$-function:
with $\xi(s)=\frac{s(s-1)}{2}\pi^{-s/2}\Gamma(s/2)\zeta(s)$ one has
$H_0(x)=\frac18\,\xi\!\bigl(\frac12+\frac{ix}2\bigr)$, so that the
Riemann hypothesis (RH) is the statement that all zeros of $H_0$ are
real. De Bruijn \cite{deBruijn} showed that if $H_t$ has only real
zeros for some $t$, the same holds for all larger $t$; Newman
\cite{Newman} showed that for some $t$ the function $H_t$ has a
non-real zero. Hence there is a constant $\LB\in\R$, the
\emph{de Bruijn--Newman constant}, such that $H_t$ has only real zeros
if and only if $t\ge\LB$. RH is equivalent to $\LB\le0$.

Newman conjectured $\LB\ge0$, which was proved by Rodgers and Tao
\cite{RodgersTao}. On the other side, a long sequence of numerical
upper bounds (de Bruijn's $\LB\le1/2$, then
\cite{CsordasNorfolkVarga,KKL} and others) led to the bound
\begin{equation}\label{eq:polymathbound}
\LB\le0.22
\end{equation}
of the Polymath~15 project \cite{Polymath}. Subsequently Platt and
Trudgian observed \cite[Corollary~2]{PlattTrudgian} that their
verification of RH to height $3\cdot10^{12}$
(Section~\ref{sec:hypi} below), combined with the second row of the
parameter table in \cite[Table~1, p.~65]{Polymath} (which requires
verified height $H>2.51\cdot10^{12}$), yields the current record
\begin{equation}\label{eq:ptbound}
\LB\le0.2 .
\end{equation}
Thus at present $0\le\LB\le0.2$. Platt and Trudgian further note
that the next entry of the same table would give $\LB<0.19$, but only
conditionally on a verified height slightly above $10^{13}$, beyond
their theorem; the bound \eqref{eq:target} below therefore
goes past both the published record \eqref{eq:ptbound} and that
conditional entry, while consuming only the established
$3\cdot10^{12}$ verification height.

\subsection{Main result}

This paper describes a claimed computer-assisted proof, due to
J.~Gomila, of the following improvement of \eqref{eq:ptbound}.

\begin{theorem}[Theorem; \textbf{unreviewed}]\label{thm:main}
\[
\LB\le\frac{893927}{5000000}=0.1787854 .
\]
\end{theorem}

The proposed proof instantiates the barrier criterion of
\cite[Theorem~1.2]{Polymath} at the exact parameters
\begin{equation}\label{eq:params}
X=6000000185827,\qquad
t_0=\frac{129}{800}=0.16125,\qquad
y_0^2=\frac{87677}{2500000}=0.0350708,
\end{equation}
with $y_0>0$, so that
\begin{equation}\label{eq:target}
t_0+\frac{y_0^2}{2}=\frac{893927}{5000000}=0.1787854 .
\end{equation}
One checks the elementary inequalities
$0<t_0<\tfrac12$, $0<y_0^2<1-2t_0=\tfrac{271}{400}$, and
$y_0^2+2t_0=\tfrac{893927}{2500000}<1$, which place the parameters in
the domain of the criterion.

\subsection{Status, and how gaps are marked}\label{sec:status}

The claimed proof combines:
\begin{enumerate}
\item two published theorem inputs --- the Polymath~15 criterion and
effective approximation \cite{Polymath}, and the Platt--Trudgian finite
verification of RH \cite{PlattTrudgian};
\item a collection of new lemmas (reproduced with proofs below),
which bridge the published estimates to the quantities computed by
the software; and
\item a large body of rigorous-interval-arithmetic certificates
(FLINT/Arb in C, plus corroborating Python interval implementations),
totalling $3{,}149{,}013$ certified finite ``window'' rows, a
standalone tail computation, and an $883$-prism barrier certificate.
\end{enumerate}
The statements quoted from the published inputs in item~(1) ---
Theorems~1.2 and~1.3 of \cite{Polymath} together with the displays
(14), (20)--(24), (71)--(72), Proposition~6.6 and Lemma~8.4 of that
paper, and Theorem~1 of \cite{PlattTrudgian} --- have been checked
verbatim against the arXiv source files of the cited papers
(1904.12438v2 and 2004.09765v1, to which all equation and theorem
numbers below refer). The one point where the current proof's usage
deliberately departs from a cited display, namely equation~(24) of
\cite{Polymath}, is discussed in Remark~\ref{rem:1044}. The
lemmas of item (2) are given complete proofs below, all verified in
the course of this review. Every computation in item (3) --- the main
finite computation (Proposition~\ref{prop:sweep}), the Dini cell
computation (Proposition~\ref{prop:dinicells}), the finite-region
error bound (Proposition~\ref{prop:finiteerr}), the tail cutoff
computation (Proposition~\ref{prop:tailcert}), the barrier-box
error bound (Proposition~\ref{prop:barriererr}), the tail-exponent
inequality (Proposition~\ref{prop:tailexp}), and the barrier prism
decomposition (Proposition~\ref{prop:prismcert}) --- has been
reproduced on an independent toolchain, by a standalone program
(supplementary directory \texttt{dan-reworking/code/}) that reads no stored data
and was reviewed line by line against the corresponding statement;
the proofs are given with the respective propositions. No
unverified steps remain flagged; the earlier framed \emph{gap /
review points} of this document have all been resolved and
removed. The verification is summarized in
Section~\ref{sec:gaps}.

% ===============================================================
\section{The Polymath criterion and effective approximation}
% ===============================================================

\subsection{The barrier criterion}

The logical skeleton of the proof is the following criterion.

\begin{theorem}[{Polymath \cite[Theorem~1.2]{Polymath}}]\label{thm:criterion}
Let $X>0$, $0<t_0<\tfrac12$ and $0<y_0^2<1-2t_0$. Suppose that:
\begin{enumerate}
\item[\textup{(i)}] \textbf{(Verified height.)} There are no zeros
$\zeta(\sigma+iT)=0$ with
\[
\frac{1+y_0}{2}\le\sigma\le1,\qquad 0\le T\le\frac X2 .
\]
\item[\textup{(ii)}] \textbf{(Final-time right region.)} There are no
zeros $H_{t_0}(x+iy)=0$ with
\[
x\ge X+\sqrt{1-y_0^2},\qquad y_0\le y\le\sqrt{1-2t_0}.
\]
\item[\textup{(iii)}] \textbf{(Intermediate-time barrier.)} There are
no zeros $H_t(x+iy)=0$ with
\[
X\le x\le X+\sqrt{1-y_0^2},\qquad
\sqrt{y_0^2+2(t_0-t)}\le y\le\sqrt{1-2t},\qquad
0\le t\le t_0 .
\]
\end{enumerate}
Then
\[
\LB\le t_0+\frac{y_0^2}{2}.
\]
\end{theorem}

Sections~\ref{sec:hypi}, \ref{sec:finite}--\ref{sec:tail}, and
\ref{sec:barrier} below describe the proof's suppliers for
hypotheses (i), (ii), and (iii) respectively, at the parameters
\eqref{eq:params}.

For hypothesis (i) one uses the normalization connecting $H_0$ to
$\zeta$: a zero $H_0(x+iy)=0$ gives
$\xi\bigl(\frac{1-y+ix}{2}\bigr)=0$; the functional equation
$\xi(s)=\xi(1-s)$ followed by complex conjugation produces the
representative zero $\xi\bigl(\frac{1+y+ix}{2}\bigr)=0$ with
$\sigma=\frac{1+y}2\ge\frac{1+y_0}2$ and ordinate $T=\frac x2$, which
is the region excluded in (i). (The proof verifies the affine
algebra of this ``sign map'' by an exact-rational script; the
functional equation and the reality symmetry
$\overline{\xi(\bar s)}=\xi(s)$ are the classical inputs.)

\subsection{The effective approximation}\label{sec:thm13}

For $t>0$ write, following \cite{Polymath},
\begin{equation}\label{eq:gt}
g_t(z)=\frac{H_t(z)}{B_t(z)},
\end{equation}
where $B_t$ is an explicit normalizing factor: in the notation of
\cite{Polymath}, $B_t(x+iy)=M_t\bigl(\frac{1+y-ix}{2}\bigr)$, where
$M_t(s)=\exp\bigl(\frac t4\alpha(s)^2\bigr)M_0(s)$ and $M_0$ is
defined in \cite[(6)]{Polymath} as the exponential of an explicit
holomorphic branch of its logarithm on $\C\setminus(-\infty,1]$
(see \cite[(7)]{Polymath}). In particular $M_0$, and hence $M_t$,
takes values in $\C\setminus\{0\}$ on its domain; since the argument
$\frac{1+y-ix}{2}$ is non-real whenever $x\neq0$, it avoids the
branch cut, and therefore $B_t(x+iy)\neq0$ for all $t\ge0$, $x>0$ and
$y>0$, as noted explicitly in \cite[\S1, after (11)]{Polymath}. This
nonvanishing is used below both in the reduction of zero-freeness of
$H_t$ to lower bounds on $|f_t|$ and in the argument-principle step
of the barrier certificate (Section~\ref{sec:barrier}).
Theorem~1.3 of \cite{Polymath} provides, in the region
\begin{equation}\label{eq:region}
x\ge200,\qquad 0\le y\le1,\qquad 0<t\le\tfrac12,
\end{equation}
an approximation
\begin{equation}\label{eq:thm13}
g_t(x+iy)=f_t(x+iy)+O_{\le}\bigl(e_A+e_B+e_{C,0}\bigr),
\end{equation}
where $O_\le(E)$ denotes a quantity of absolute value at most $E$, the
main term is the finite sum (equation (14) of \cite{Polymath})
\begin{equation}\label{eq:ft}
f_t(x+iy)
=\sum_{n=1}^{N}b_t(n)\,n^{-s_*}
+\gamma\sum_{n=1}^{N}n^{y}\,b_t(n)\,n^{-\overline{s_*}-\kappa},
\qquad
b_t(n)=\exp\Bigl(\frac t4\log^2 n\Bigr),
\end{equation}
(we write $b_t(n)$ for the quantity denoted $b_n^t$ in
\cite{Polymath}; all other symbols in \eqref{eq:ft} agree with the
notation there, including the placement of $\kappa$ alongside the
conjugated exponent $\overline{s_*}$ in the second sum)
with:
\begin{itemize}
\item the window index
$N=\bigl\lfloor\sqrt{\tfrac{x}{4\pi}+\tfrac{t}{16}}\bigr\rfloor$;
\item explicitly defined quantities $\gamma=\gamma(x,y,t)$,
$s_*=s_*(x,y,t)$, $\kappa=\kappa(x,y,t)$ satisfying (by equations
(20)--(22) of \cite{Polymath}) the bounds
\begin{align}
|\gamma|&\ \le\ e^{0.02y}\Bigl(\frac{x}{4\pi}\Bigr)^{-y/2},
\label{eq:gammabound}\\
\Rea s_*&\ \ge\ \frac{1+y}{2}+\frac t4\log\frac{x}{4\pi}
-\frac{t}{2x^2}\Bigl(1-3y+\frac{4y(1+y)}{x^2}\Bigr)_{\!+},
\label{eq:resbound}\\
|\kappa|&\ \le\ \frac{ty}{2(x-6)};
\label{eq:kappabound}
\end{align}
\item and error terms $e_A,e_B,e_{C,0}$ with explicit majorants
(equations (71)--(72) and (23)--(24) of \cite{Polymath}).
\end{itemize}
Since $B_t\ne0$, nonvanishing of $H_t$ on a region follows from the
strict inequality $|f_t|>e_A+e_B+e_{C,0}$ there; this implication is
itself stated in \cite{Polymath} as Corollary~1.4 (``Criterion for
non-vanishing'').

\subsection{The conservative error constant $10.50$}\label{sec:weld}

The current proof deliberately does \emph{not} use the constant $10.44$
displayed in equation~(24) of \cite{Polymath} for the majorization of
$e_{C,0}$.
Instead it derives a
slightly weaker bound directly from Proposition~6.6(vi) of
\cite{Polymath}, which contains the expression
\[
\exp\!\left(\frac{3\bigl|\log\frac{x}{4\pi}+\frac{i\pi}2\bigr|+3.58}{x-8.52}\right)
\left(1+\frac{1.24\,(3^y+3^{-y})}{N-0.125}+\frac{6.92}{x-12}\right).
\]
On the theorem domain $x\ge200$ one has $x-8.52>x-12>0$; applying
$1+u\le e^u$ to the second factor and enlarging the denominator
$x-8.52$ to $x-12$ yields the \emph{conservative corollary}
\begin{equation}\label{eq:ec0}
e_{C,0}\le
\Bigl(\frac{x}{4\pi}\Bigr)^{-(1+y)/4}
\exp\!\left(
-\frac{t}{16}\log^2\frac{x}{4\pi}
+\frac{1.24\,(3^y+3^{-y})}{N-0.125}
+\frac{3\bigl|\log\frac{x}{4\pi}+\frac{i\pi}2\bigr|+10.50}{x-12}
\right),
\end{equation}
with $3.58+6.92=\tfrac{21}2=10.50$ exactly. Every numerical error
computation below (finite range, tail, and barrier) uses
\eqref{eq:ec0}, and uses the denominator $x-6.66$, which appears in
equation (23) and in Proposition~6.6(iv),(v) of \cite{Polymath}, for
$e_A+e_B$.

\begin{remark}\label{rem:1044}
Equation (24) of \cite{Polymath} displays the bound \eqref{eq:ec0}
with the constant $10.44$ in place of $10.50$. The sentence following
Proposition~6.6 in \cite{Polymath} states that (24) is obtained from
part (vi) of the proposition by ``the inequality $1+u\le\exp(u)$''
followed by the bound $\frac{1}{x-8.52}\le\frac{1}{x-12}$; carrying
out this recipe yields the constant $3.58+6.92=10.50$, and we are not
aware of an argument in \cite{Polymath} producing the displayed
$10.44$. Accordingly, the current proof does not rely on (24): all
error computations use only
\eqref{eq:ec0}, which follows rigorously from Proposition~6.6(vi) by
the two elementary steps just described, uniformly on the domain
\eqref{eq:region}. (The weaker constant costs a factor
$e^{0.06/(x-12)}$, which is utterly negligible at
$x\approx6\times10^{12}$.)
\end{remark}

% ===============================================================
\section{Hypothesis (i): the verified height}\label{sec:hypi}
% ===============================================================

Platt and Trudgian \cite[Theorem~1]{PlattTrudgian} proved that every
nontrivial zero of $\zeta$ with ordinate in $(0,T_{\rm PT}]$,
\[
T_{\rm PT}=3000175332800,
\]
lies on the critical line. (Their statement reads: ``The Riemann
hypothesis is true up to height $3\,000\,175\,332\,800$. That is, the
lowest $12\,363\,153\,437\,138$ non-trivial zeroes $\rho$ have
$\Re\rho=1/2$.'' The method --- counting sign changes of the completed
zeta function on the critical line and confirming by Turing's method
that all zeros are accounted for --- yields the off-line exclusion
form used here.) Hypothesis (i) of
Theorem~\ref{thm:criterion} requires the absence of zeros with
$\sigma\ge\frac{1+y_0}2>\frac12$ and $0\le T\le X/2$. Since
\[
\frac X2=3000000092913.5,
\qquad
T_{\rm PT}-\frac X2=\frac{350479773}{2}>0,
\]
the verified region covers all ordinates $0<T\le X/2$ with margin
$\approx1.75\times10^8$. On the segment $T=0$, $\tfrac12<\sigma<1$,
one argues directly: the Dirichlet eta function satisfies
$\eta(\sigma)=\sum_{n\ge1}(-1)^{n-1}n^{-\sigma}>0$ for
$0<\sigma<1$ (group consecutive terms), while
$\eta(\sigma)=(1-2^{1-\sigma})\zeta(\sigma)$ with
$1-2^{1-\sigma}<0$; hence $\zeta(\sigma)<0\ne0$ there. At $\sigma=1$,
$\zeta$ has a pole. This supplies hypothesis (i).

As with any use of a computationally established theorem, this step
inherits the correctness of the Platt--Trudgian computation itself (a
rigorous interval-arithmetic computation, independent of anything in
the present paper); we accept it here as a published-theorem
input.

% ===============================================================
\section{Hypothesis (ii), finite range: the window scan}\label{sec:finite}
% ===============================================================

Hypothesis (ii) concerns nonvanishing of $H_{t_0}$ on
\begin{equation}\label{eq:hypiiregion}
x\ge x_*:=X+\sqrt{1-y_0^2},
\qquad
y_0\le y\le\sqrt{1-2t_0}.
\end{equation}
Gomila's proof splits this region at the abscissa corresponding to
the window index $N_*=3840000$: the finite range
$690988\le N\le3840000$ is handled in this section by
$3{,}149{,}013$ individually certified interval computations, and
the infinite range $N\ge3840000$ in Section~\ref{sec:tail} by a
single analytic contraction argument. The two ranges overlap on a
full window, eliminating endpoint risk.

\subsection{Windows and the freeze lemma}\label{sec:freeze}

At $t=t_0$ define
\[
x_N=4\pi\Bigl(N^2-\frac{t_0}{16}\Bigr),
\qquad
W_N=[x_N,x_{N+1}),
\]
so that $x\in W_N$ iff the window index of $x$ equals $N$. Exact
rational Machin-type bounds for $\pi$ and exact integer square-root
brackets (independently confirmed by a 400-bit interval computation)
give
\begin{equation}\label{eq:xstarwindow}
x_{690988}<x_*<x_{690989},
\end{equation}
with margins exceeding $5.37\times10^6$ and $1.19\times10^7$
respectively; thus $x_*$ lies strictly inside $W_{690988}$, and the
half-open windows tile $[x_*,\,x_{3840001})$ with no gaps or overlaps.

The certified computations evaluate, for each $N$, the majorants of
\eqref{eq:gammabound}--\eqref{eq:kappabound} \emph{frozen at the
left edge} $x=x_N$ of the window.
The following elementary lemma justifies this.

\begin{lemma}[Window freeze]\label{lem:freeze}
For $y_0\le y\le\sqrt{1-2t_0}$ define
\begin{align*}
G(x,y) &= e^{y/50}\Bigl(\frac{x}{4\pi}\Bigr)^{-y/2},
\\
K(x,y) &= \frac{t_0y}{2(x-6)},
\\
\Sigma(x,y) &= \frac{1+y}{2}+\frac{t_0}{4}\log\frac{x}{4\pi}
-\frac{t_0}{2x^2}\Bigl(1-3y+\frac{4y(1+y)}{x^2}\Bigr)_{\!+}.
\end{align*}
Then for every $N$ in the finite range, every $x\in W_N$ and every $y$
in the stated range,
\[
G(x,y)\le G(x_N,y),\qquad
K(x,y)\le K(x_N,y),\qquad
\Sigma(x,y)\ge\Sigma(x_N,y).
\]
\end{lemma}

\begin{proof}
$\partial_xG=-\frac{y}{2x}G<0$ and
$\partial_xK=-\frac{t_0y}{2(x-6)^2}<0$ give the first two
inequalities. For the third, write
$h(x,y)=1-3y+A(y)/x^2$ with $A(y)=4y(1+y)\ge0$; then
$h_x=-2A/x^3\le0$, so $h_+$ is nonincreasing in $x$, and since
$x^{-2}$ is positive and decreasing, the product $x^{-2}h_+$ is
nonincreasing. The logarithmic term is strictly increasing, so
$\Sigma$ is strictly increasing in $x$, including across the kink
$h=0$ (both one-sided derivatives there are positive).
\end{proof}

Note that $e^{y/50}=e^{0.02y}$, matching \eqref{eq:gammabound}; $G,K,\Sigma$
frozen at $x_N$ are precisely the quantities $g_N$ (bound for
$|\gamma|$), $k_N$ (bound for $|\kappa|$) and $\sigma_N$ (lower bound
for $\Rea s_*$) used in the definition of the certified quantities of Section~\ref{sec:stored}.

\subsection{The mollified lower bound for $|f_{t_0}|$}\label{sec:native}

Fix a window $N$, a height $y>0$, and a finite set of auxiliary primes
$\calP$ with $D=\prod_{p\in\calP}p$. Let $g_N,k_N,\sigma_N$ be as
above, and introduce the real \emph{Euler mollifier}
\[
E(s_*)=\prod_{p\in\calP}\bigl(1-b_{t_0}(p)\,p^{-s_*}\bigr)
=\sum_{d\mid D}\lambda_d\,d^{-s_*},
\qquad
\lambda_d=(-1)^{\omega(d)}\prod_{p\mid d}b_{t_0}(p)\in\R .
\]
Define the convolved coefficients (for $1\le n\le DN$)
\[
B_{N,n}=\sum_{\substack{d\mid D,\ d\mid n\\ n/d\le N}}
\lambda_d\,b_{t_0}(n/d),
\qquad
A_{N,n}=\sum_{\substack{d\mid D,\ d\mid n\\ n/d\le N}}
\lambda_d\,b_{t_0}(n/d)\,(n/d)^y,
\]
so $B_{N,1}=A_{N,1}=1$, and set
\[
M_N=\sum_{d\mid D}|\lambda_d|\,d^{-\sigma_N},
\qquad
C_N=\sum_{m=2}^{N}b_{t_0}(m)\bigl(m^{k_N}-1\bigr)m^{y-\sigma_N},
\]
\begin{equation}\label{eq:LN}
L_N=\frac{1-g_N-\sum_{n=2}^{DN}
\bigl(|B_{N,n}|+g_N|A_{N,n}|\bigr)n^{-\sigma_N}}{M_N}
\;-\;g_NC_N .
\end{equation}

\begin{lemma}[Mollified lower bound]\label{lem:native}
If $L_N>0$, then $|f_{t_0}(x+iy)|\ge L_N$ for all $x\in W_N$.
\end{lemma}

\begin{proof}
Write \eqref{eq:ft} as $f_t=A+\gamma C_\kappa$ with
$A=\sum_{m\le N}b_t(m)m^{-s_*}$ and
$C_\kappa=\sum_{m\le N}m^yb_t(m)m^{-\overline{s_*}-\kappa}$. Exact
Dirichlet convolution gives the two identities
\begin{equation}\label{eq:conv}
E(s_*)\,A=\sum_{n=1}^{DN}B_{N,n}\,n^{-s_*},
\qquad
\overline{E(s_*)}\;C_0=\sum_{n=1}^{DN}A_{N,n}\,n^{-\overline{s_*}} ,
\end{equation}
where the second identity uses that the $\lambda_d$ are \emph{real}
(so that $\overline{E(s_*)}=\sum_d\lambda_d d^{-\overline{s_*}}$).
By the triangle inequality and $\Rea s_*\ge\sigma_N$,
\[
|EA|\ge1-\sum_{n=2}^{DN}|B_{N,n}|n^{-\sigma_N},
\qquad
|\overline EC_0|\le1+\sum_{n=2}^{DN}|A_{N,n}|n^{-\sigma_N}.
\]
For complex $\kappa$ with $|\kappa|\le k_N$,
\begin{equation}\label{eq:kappaest}
|m^{-\kappa}-1|=\bigl|e^{-\kappa\log m}-1\bigr|
\le e^{|\kappa|\log m}-1\le m^{k_N}-1,
\end{equation}
whence $|\overline E\,(C_\kappa-C_0)|\le|E|\,C_N$. Combining, and
using $|\gamma|\le g_N$,
\[
|E|\,|f_t|\;\ge\;Q_N-g_N|E|C_N,
\qquad
Q_N:=1-g_N-\sum_{n=2}^{DN}\bigl(|B_{N,n}|+g_N|A_{N,n}|\bigr)n^{-\sigma_N}.
\]
If $L_N>0$ then $Q_N=M_N(L_N+g_NC_N)>0$; the displayed inequality then
forces $E\ne0$ (otherwise $0\ge Q_N>0$). Dividing by $|E|$ and using
$|E|\le M_N$ together with $Q_N>0$,
\[
|f_t|\ \ge\ \frac{Q_N}{|E|}-g_NC_N
\ \ge\ \frac{Q_N}{M_N}-g_NC_N=L_N. \qedhere
\]
\end{proof}

The point of Lemma~\ref{lem:native} is that any certified
per-window lower bound for $L_N$ is \emph{already} in the same
normalized units as the additive error in \eqref{eq:thm13}; no
further Euler-factor conversion intervenes.

\subsection{The certified lower bounds}\label{sec:stored}

For each integer $N$ with $690988\le N\le 3840000$, define a set of
auxiliary primes
\begin{equation}\label{eq:PN}
\calP(N)=
\begin{cases}
\{2,3,5,7,11\}, & 690988\le N\le728999,\\
\{2,3,5,7\}, & 729000\le N\le818999,\\
\{2,3,5\}, & 819000\le N\le1027999,\\
\{2,3\}, & 1028000\le N\le3840000,
\end{cases}
\end{equation}
and let $L_N$ denote the quantity \eqref{eq:LN} formed with the
prime set $\calP(N)$, at $t=t_0$ and $y=y_0$, with the constants
$g_N=G(x_N,y_0)$, $k_N=K(x_N,y_0)$, $\sigma_N=\Sigma(x_N,y_0)$
frozen at the window's left edge as in Lemma~\ref{lem:freeze}. (The
number of terms in $L_N$ grows with $|\calP(N)|$; the choice
\eqref{eq:PN} uses the larger, more effective mollifiers only on the
ranges of $N$ where they are needed.)

\begin{proposition}[Certified lower bounds for $L_N$]\label{prop:sweep}
For every integer $690988\le N\le3840000$,
\[
L_N\ \ge\ T_{\min}\ =\ 791366\times10^{-12}\ >\ 0 .
\]
More precisely, the minimum of $L_N$ over each of the four ranges in
\eqref{eq:PN} is bounded below as listed in
Table~\ref{tab:primesets}; the global bound $T_{\min}$ is attained
at $N=690988$.
\end{proposition}

\begin{table}[ht!]
\centering
\begin{tabular}{@{}llr@{}}
\toprule
$\calP(N)$ & range of $N$ & certified lower bound for $\min L_N$\\
\midrule
$\{2,3,5,7,11\}$ & $690988\le N\le728999$ & $0.000000791366$\\
$\{2,3,5,7\}$    & $729000\le N\le818999$ & $0.000315112481$\\
$\{2,3,5\}$      & $819000\le N\le1027999$ & $0.000305788816$\\
$\{2,3\}$        & $1028000\le N\le3840000$ & $0.000309285478$\\
\bottomrule
\end{tabular}
\caption{The auxiliary prime sets \eqref{eq:PN} and, for each of the
four ranges of $N$, the certified lower bound for $\min_N L_N$
established by the verification run described in the proof of
Proposition~\ref{prop:sweep}.}
\label{tab:primesets}
\end{table}

\begin{proof}
The proof is computer-assisted; 
the verification program \texttt{prop43\_proof.c}, included with
this manuscript as supplementary material (directory
\texttt{code/prop43}), computes for each $N$ an outward-rounded
interval enclosure of a quantity provably at most $L_N$, and prints
the floor of $10^{12}$ times the lower endpoint of that enclosure.
Each printed value is therefore a certified lower bound for $L_N$.
The program is derived from Gomila's source (as documented in its
header) and was reviewed line by line as part of this audit. It has
two evaluation modes: a direct term-by-term interval evaluation of
\eqref{eq:LN}, and a production sweep that majorizes the sum in
\eqref{eq:LN} by activation-frozen weights (conservative because the
factor $g$ decreases in $N$) and evaluates the result through
certified Taylor moments with an explicit Lagrange remainder. The
review verified that every step of both modes is conservatively
directed, and the two modes produce the identical bound at the
extremal index $N=690988$.

The program was compiled against an independently installed
FLINT/Arb~3.6.0 and run at 256-bit precision over the four
configurations of \eqref{eq:PN}, in each case at $y=y_0$ and over a
closed rational $t$-interval containing $t_0$ (for the first range,
the singleton $\{129/800\}$). Every one of the $3{,}149{,}013$
values of $N$ produced a certified positive bound, with no
uncertified rows; the concatenated outputs cover
$690988\le N\le3840000$ with no gaps or duplicates; and the
certified per-range minima are as in Table~\ref{tab:primesets}, with
global minimum $791366\times10^{-12}$ at $N=690988$. Per-range
summaries (\texttt{runs/*\_summary.txt}) and the complete per-$N$
records (\texttt{runs/*\_sweep.out}) are included in the
supplementary material. 
\end{proof}

\subsection{Transfer to the full height range}\label{sec:dini}

Proposition~\ref{prop:sweep} concerns the single height $y=y_0$,
while hypothesis (ii) requires nonvanishing for all
$y\in[y_0,\sqrt{1-2t_0}]$. For fixed $N$ write $L_N(y)$ for the
quantity \eqref{eq:LN} with its explicit $y$-dependence restored
(the constants $g_N,k_N,\sigma_N$ frozen in $x$ as before, but
evaluated at height $y$; likewise the coefficients $A_{N,n}$ and the
sums $M_N,C_N$), so that $L_N=L_N(y_0)$. Throughout this subsection
$t=t_0$ and $N$ are fixed, and we abbreviate
\[
q_N=N^2-\frac{t_0}{16},
\qquad
g=\frac1{50}-\frac12\log q_N\ (<0),
\qquad
\rho_N=\frac{t_0}{2(x_N-6)},
\]
so that $g_N(y)=e^{yg}$ exactly and $k_N(y)=\rho_Ny$. The bridge
between Proposition~\ref{prop:sweep} and the full height range is
the following statement.

\begin{proposition}[Height transfer]\label{prop:transfer}
For every integer $690988\le N\le3840000$ and every
$y\in[y_0,\sqrt{1-2t_0}]$,
\[
L_N(y)\ \ge\ L_N(y_0).
\]
\end{proposition}

The proof establishes three monotonicity statements: as $y$
increases over $[y_0,\sqrt{1-2t_0}]$, (a) the majorant sum
\[
\mathcal F_{N,\calP}(y)
=g_N(y)+\sum_{n=2}^{DN}
\bigl(|B_{N,n}|+g_N(y)|A_{N,n}(y)|\bigr)n^{-\sigma_N(y)}
\]
subtracted in \eqref{eq:LN} is nonincreasing; (b) the mollifier
bound $M_N(y)$ is positive and nonincreasing; and (c) the correction
term $g_N(y)C_N(y)\ge0$ is nonincreasing. Statement (a) is where an
earlier, now-abandoned shortcut of Gomila's failed: the crude
majorant obtained by discarding the negative term $g_N g|A|$ below
is not monotone, and the proof must retain the cancellation between
the $A$-side terms. We build it up in three lemmas and one certified
computation.

\begin{lemma}[Slope of the frozen exponent]\label{lem:sigmaslope}
For $y_0\le y_1\le y_2\le\sqrt{1-2t_0}$,
\[
\sigma_N(y_2)-\sigma_N(y_1)\ \ge\ \tfrac12\,(y_2-y_1).
\]
\end{lemma}

\begin{proof}
$\sigma_N(y)=\frac{1+y}2+\frac{t_0}4\log q_N
-\frac{t_0}{2x_N^2}\bigl(1-3y+\frac{4y(1+y)}{x_N^2}\bigr)_+$. The
inner expression has $y$-derivative $-3+\frac{4(1+2y)}{x_N^2}<0$
(as $4(1+2y)\le12<3x_N^2$), so its positive part is nonincreasing
in $y$ and the subtracted term is nondecreasing; the explicit term
$\frac{1+y}2$ contributes slope $\frac12$.
\end{proof}

\begin{lemma}[Monotone normalizer and correction]\label{lem:bc}
On $[y_0,\sqrt{1-2t_0}]$: \textup{(b)} $M_N(y)$ is positive and
nonincreasing; \textup{(c)} $g_N(y)C_N(y)$ is nonnegative and
nonincreasing.
\end{lemma}

\begin{proof}
(b) Each term $|\lambda_d|d^{-\sigma_N(y)}$ is positive; for
$d\ge2$ it is nonincreasing by Lemma~\ref{lem:sigmaslope}, and the
$d=1$ term is constant.

(c) Nonnegativity is clear ($m^{k_N(y)}\ge1$). Fix
$y_0\le y_1<y_2\le\sqrt{1-2t_0}$ and put $\Delta=y_2-y_1$. First,
$g_N(y_2)=g_N(y_1)e^{g\Delta}$ exactly. Next, for each
$2\le m\le N$, the term
$b_t(m)\,(m^{\rho_Ny}-1)\,m^{\,y-\sigma_N(y)}$ of $C_N$ grows by at
most the factor $e^{(\frac12\log m+\rho_N\log m+\frac1{y_0})\Delta}$:
the exponent $y-\sigma_N(y)$ grows by at most $\frac12\Delta$
(Lemma~\ref{lem:sigmaslope}); and by the mean value theorem applied
to $y\mapsto\log(m^{\rho_Ny}-1)$, whose derivative is
$\rho_N\log m\,\frac{e^z}{e^z-1}$ with $z=\rho_Ny\log m>0$, the
bound $\frac{e^z}{e^z-1}\le1+\frac1z$ gives a growth rate at most
$\rho_N\log m+\frac1{y}\le\rho_N\log m+\frac1{y_0}$. Since a sum of
positive terms grows no faster than its fastest-growing term, and
$m\le N$,
\[
g_N(y_2)C_N(y_2)\ \le\
g_N(y_1)C_N(y_1)\,
\exp\Bigl[\Bigl(g+\tfrac12\log N+\rho_N\log N
+\tfrac1{y_0}\Bigr)\Delta\Bigr].
\]
It remains to check that the rate is negative for every $N$ in the
range. First, $g+\frac12\log N=\frac1{50}-\frac12\log\frac{q_N}N$
and $\frac{q_N}N=N-\frac{t_0}{16N}>N-1\ge690987>e^{13.44}$ (as
$e^{13.44}<442414\times1.5528<687000$), so
$g+\frac12\log N<0.02-6.72$. Second,
$\frac1{y_0}<\frac{10^7}{1872719}<5.3399$, since
$y_0^2=\frac{350708}{10^7}>\bigl(\frac{1872719}{10^7}\bigr)^2
=0.03507076\ldots$ gives $y_0>\frac{1872719}{10^7}$. Third,
uniformly over the range,
$\rho_N\log N\le\frac{t_0\log(3.84\times10^6)}{2(x_{690988}-6)}
<\frac{0.17\times16}{10^{13}}<10^{-12}$. Hence the rate is at most
$0.02+5.3399+10^{-12}-6.72<-1.36<0$.
\end{proof}

\begin{lemma}[Dini bound in divisor-cell form]\label{lem:dini}
Let $2\le n\le DN$, $L=\log n$, and let
$\mathcal D_{N,n}=\{d\mid D:\ d\mid n,\ n\le dN\}$ be the active
divisor set (which does not depend on $y$), assumed nonempty. With
$\ell_d=\log d$ and
\begin{align*}
r_d(L) &= \exp\Bigl\{\frac{t_0}4\Bigl(\sum_{p\mid d}\log^2p
+\ell_d^2-2L\ell_d\Bigr)\Bigr\},
\\
C_0 &= \sum_{d\in\mathcal D_{N,n}}(-1)^{\omega(d)}r_d(L),
\\
C(y) &= \sum_{d\in\mathcal D_{N,n}}(-1)^{\omega(d)}r_d(L)\,d^{-y},
\\
C'(y) &= -\sum_{d\in\mathcal D_{N,n}}(-1)^{\omega(d)}r_d(L)\,
d^{-y}\ell_d,
\end{align*}
one has $B_{N,n}=b_t(n)\,C_0$,
$A_{N,n}(y)=b_t(n)\,e^{yL}C(y)$ and
$\frac{d}{dy}A_{N,n}(y)=b_t(n)\,e^{yL}(LC(y)+C'(y))$, and the upper
right Dini derivative of the corresponding summand of
$\mathcal F_{N,\calP}$ satisfies
\[
D^+_y\Bigl[\bigl(|B_{N,n}|+g_N(y)|A_{N,n}(y)|\bigr)
n^{-\sigma_N(y)}\Bigr]
\ \le\
b_t(n)\,n^{-\sigma_N(y)}
\Bigl(G\,h-\frac L2\,|C_0|\Bigr),
\]
where $G=G(y)=e^{y(g+L)}$ and
$h=h(y)=|LC(y)+C'(y)|+\bigl(g-\tfrac L2\bigr)|C(y)|$.
\end{lemma}

\begin{proof}
The identities follow from
$\lambda_d\,b_t(n/d)=(-1)^{\omega(d)}b_t(n)\,r_d(L)$ --- indeed
$\log b_t(n/d)=\frac{t_0}4(L-\ell_d)^2
=\log b_t(n)+\frac{t_0}4(\ell_d^2-2L\ell_d)$ and
$\lambda_d=(-1)^{\omega(d)}\exp(\frac{t_0}4\sum_{p\mid d}\log^2p)$
--- together with $(n/d)^y=e^{yL}d^{-y}$; note also
$g_N(y)e^{yL}=e^{y(g+L)}=G$. Since $A=A_{N,n}(y)$ is real-analytic
in $y$,
\[
D^+_y\bigl\{g_N|A|\bigr\}
=\begin{cases}
g_N\bigl(g|A|+\operatorname{sgn}(A)A'\bigr),&A\neq0,\\
g_N|A'|,&A=0,
\end{cases}
\ \le\ g_N\bigl(g|A|+|A'|\bigr)
\]
in both cases. Both factors of the summand are nonnegative and
locally Lipschitz, $|B|$ is constant in $y$, and by
Lemma~\ref{lem:sigmaslope},
$D^+_y\,n^{-\sigma_N(y)}\le-\frac L2\,n^{-\sigma_N(y)}$; the
product rule for upper Dini derivatives of nonnegative Lipschitz
functions then gives
\[
D^+_y\bigl[(|B|+g_N|A|)n^{-\sigma_N}\bigr]
\le n^{-\sigma_N}\Bigl(g_N\bigl(g|A|+|A'|\bigr)
-\frac L2\bigl(|B|+g_N|A|\bigr)\Bigr).
\]
Substituting the identities and factoring out $b_t(n)>0$ yields the
display.
\end{proof}

\begin{lemma}[Cell decomposition and $q$-freezing]\label{lem:cells}
\textup{(i)} For $2\le n\le DN$ put $G_0=\gcd(n,D)$ and list the
divisors of $G_0$ as $1=d_0<d_1<\dots<d_s$ (so $s+1=\tau(G_0)$),
with $d_{-1}=0$. On the ratio sector $d_{j-1}<n/N\le d_j$
($0\le j\le s$) the active set is exactly
$\mathcal D_{N,n}=\{d_j,\dots,d_s\}$, and for $n/N>d_s$ it is
empty (the summand vanishes identically). The number of pairs
(mask, sector) with nonempty active set is
$\sum_{G_0\mid D}\tau(G_0)=3^{|\calP|}$.

\textup{(ii)} With $c=|C(y)|$ and
$K=|LC(y)+C'(y)|-\frac L2c$, so that
$G\,h_+=U(g):=e^{y(L+g)}\,[K+gc]_+$, the function
$\tilde g\mapsto U(\tilde g)$ is nondecreasing.

\textup{(iii)} On the sector of \textup{(i)}, within the range
$N_{\min}\le N\le N_{\max}$ of the prime set,
$q_N\ge\max\bigl(N_{\min}^2,\,e^{2L}/d_j^2\bigr)-\frac{t_0}{16}$;
since $g$ is decreasing in $q_N$, replacing $g$ by its value at
this smallest feasible $q_N$ can only increase $U(g)$.
\end{lemma}

\begin{proof}
(i) $d\in\mathcal D_{N,n}$ iff $d\mid G_0$ and $n\le dN$, i.e.\
$d\ge n/N$; on the stated sector these are exactly
$d_j,\dots,d_s$. Summing $\tau(G_0)$ over the divisors $G_0$ of the
squarefree $D$ gives $\prod_{p\in\calP}(\tau(1)+\tau(p))
=3^{|\calP|}$. (ii) Where $K+\tilde gc>0$,
$U'(\tilde g)=e^{y(L+\tilde g)}\bigl(y(K+\tilde gc)+c\bigr)\ge0$;
where $K+\tilde gc\le0$, $U=0$; and $U$ is continuous. (iii)
$n\le d_jN$ gives $N\ge e^L/d_j$, and $N\ge N_{\min}$; $q_N$ is
increasing in $N$ and $g$ decreasing in $q_N$.
\end{proof}

By Lemma~\ref{lem:cells}, for each of the four prime sets the
quantity $G\,h_+$ is bounded by an expression depending only on the
cell, on $L$ and on $y$; the following finite family of
inequalities is therefore checkable by a finite computation.

\begin{proposition}[Certified cell inequalities]\label{prop:dinicells}
For each of the four prime sets of \eqref{eq:PN}, every $N$ in the
corresponding range, every $2\le n\le DN$ with
$\mathcal D_{N,n}\neq\emptyset$, and every
$y\in[y_0,\sqrt{1-2t_0}]$, in the notation of
Lemma~\ref{lem:dini},
\[
G\,h_+\ \le\ \frac L2\,|C_0| .
\]
\end{proposition}

\begin{proof}
The proof is computer-assisted; the verification program
\texttt{prop49\_proof.c}, included with this manuscript as
supplementary material (directory \texttt{code/prop49}; derived
from Gomila's source as documented in its header and reviewed line
by line as part of this audit), reads no stored data and fails
closed at every decision.

For each prime set, the program enumerates the $3^{|\calP|}$
(mask, ratio-sector) cells of Lemma~\ref{lem:cells}(i): for each
divisor mask $G_0\mid D$ and each divisor $d_j$ of $G_0$, the
active set $\{d_j,\dots,d_s\}$ with the sector constraint
$d_{j-1}<n/N\le d_j$. The $(L,y)$-domain of the cell, namely
$$L\in[\max(\log2,\log(N_{\min}d_{j-1})),\log(N_{\max}d_j)]$$ and
$y\in[y_0,\sqrt{1-2t_0}]$, is covered by closed binary64
rectangles whose endpoints are padded outward by $2\times10^{-12}$;
the containment of each exact logarithm in its padded double
endpoint, and of the exact $y_0$ and $\sqrt{1-2t_0}$ in the fixed
$y$-box, is itself certified by ball arithmetic before use. On
each rectangle the program freezes $q_N$ at the smallest feasible
value $\max\bigl(N_{\min}^2,e^{2L}/d_j^2\bigr)-\frac{t_0}{16}$
(splitting the $L$-range at the transition point with overlapping
padding; conservative by Lemma~\ref{lem:cells}(ii)--(iii)), forms
ball enclosures of $C_0$, $C(y)$, $C'(y)$ from the exact identity
$r_d(L)=e^{\frac{t_0}4(\sum_{p\mid d}\log^2p+\ell_d^2)}
e^{-\frac{t_0}2L\ell_d}$, and certifies either
$h\le0$ or the strict ball inequality $G\,h_+<\frac L2|C_0|$,
bisecting the longer side adaptively (depth at most $42$) and
failing closed on any unresolved leaf. Before starting, the program also verifies numerically, at the
smallest window, the two elementary inequalities
$4(1+2y)<3x_{690988}^2$ and $g<0$: the first is the assumption
underlying Lemma~\ref{lem:sigmaslope}, and the second makes the
standalone term $g_N(y)=e^{yg}$ decreasing in the proof of
Proposition~\ref{prop:transfer}. Both have already been verified
analytically in the proofs just cited, and the numerical check adds
nothing to the present proof; it is a deliberate redundancy of
Gomila's program, which refuses to run when these inequalities
fail, so that a passing run remains meaningful even if the
parameters hard-coded in the program were ever modified.

The fresh runs, at both $180$- and $256$-bit precision, certify
all four prime sets with no unresolved leaves ($297490$, $52239$,
$9290$ and $237$ leaves respectively, identical covers at both
precisions), and the aggregated certified ratio bound
\[
\sup\ \frac{G\,h_+}{\frac L2|C_0|}\ \le\ 0.99999860767275095
\ <\ 0.9999988\ <\ 1
\]
(the closing comparison against the fixed threshold $0.9999988$
was added in this review; Gomila's original certified only the
per-leaf inequalities and printed the worst ratio without checking
it). The worst leaf lies in the five-prime range at
$L\approx1.042$--$1.054$, $y\approx0.1997$--$0.2022$; its margin
$1-\text{ratio}\approx1.4\times10^{-6}$ is the second-thinnest in
the entire proof. Each run completes in under ten seconds; the
outputs are included as
\texttt{code/prop49/runs/prop49\_report\_180.txt}
and \texttt{prop49\_report\_256.txt}.
\end{proof}

\begin{proof}[Proof of Proposition~\ref{prop:transfer}]
\emph{Statement (a).} For each $n$ with nonempty active set,
Lemma~\ref{lem:dini} and Proposition~\ref{prop:dinicells} give
$D^+_y\le b_t(n)n^{-\sigma_N}(G\,h-\frac L2|C_0|)\le0$ on
$[y_0,\sqrt{1-2t_0}]$ (if $h\le0$ both terms are nonpositive;
otherwise apply the proposition). Summands with empty active set
vanish identically, and the standalone term $g_N(y)=e^{yg}$ is
decreasing. Each summand is continuous, and a continuous function
with nonpositive upper right Dini derivative on an interval is
nonincreasing (if $f(y_2)>f(y_1)$, put
$\eps=\frac{f(y_2)-f(y_1)}{2(y_2-y_1)}$ and let $y^*$ be the
supremum of the set where $f(y)\le f(y_1)+\eps(y-y_1)$; then
$y^*<y_2$ and $D^+f(y^*)\ge\eps>0$). Hence
$\mathcal F_{N,\calP}$ is nonincreasing.

\emph{Conclusion.} By Proposition~\ref{prop:sweep},
$1-\mathcal F_{N,\calP}(y_0)=M_N(y_0)L_N(y_0)+M_N(y_0)\,
g_N(y_0)C_N(y_0)\ge M_N(y_0)\,T_{\min}>0$. By (a) the numerator
$1-\mathcal F_{N,\calP}(y)$ is nondecreasing, hence positive on the
whole interval; by Lemma~\ref{lem:bc}(b) the denominator $M_N(y)$
is positive and nonincreasing; so the quotient
$(1-\mathcal F_{N,\calP}(y))/M_N(y)$ is nondecreasing. Subtracting
the nonnegative nonincreasing correction $g_N(y)C_N(y)$
(Lemma~\ref{lem:bc}(c)) preserves this, and
$L_N(y)\ge L_N(y_0)$ follows.
\end{proof}

\subsection{Error bounds and the finite-range conclusion}
\label{sec:finiteerr}

The remaining ingredient is a uniform bound on the errors of
Theorem~1.3 over the finite region.

\begin{proposition}[Uniform error bound, finite region]\label{prop:finiteerr}
For all $x\ge x_*$ and $y_0\le y\le\sqrt{1-2t_0}$, the error terms
of Theorem~1.3 at $t=t_0$, bounded using equation (23) of
\cite{Polymath} for $e_A+e_B$ and the conservative corollary
\eqref{eq:ec0} for $e_{C,0}$, satisfy
\[
e_A+e_B\le2.06\times10^{-12},\qquad
e_{C,0}\le0.000000233492848188649183,
\]
and
\[
e_A+e_B+e_{C,0}\ \le\ E_{\max}:=0.000000233494905212337849 .
\]
(The third bound is certified directly and is sharper than the sum
of the two displayed bounds, which are rounded upward for display.)
\end{proposition}

\begin{proof}
The proof is computer-assisted; the verification program
\texttt{prop410\_proof.py}, included with this manuscript as
supplementary material (directory \texttt{code/prop410}; derived
from Gomila's source as documented in its header and reviewed line
by line as part of this audit), computes outward-rounded interval
enclosures ($220$-bit precision) of majorants of the right-hand
sides of equation (23) of \cite{Polymath} and of the corollary
\eqref{eq:ec0}. Its only inputs are the exact rational parameters
$t_0$, $y_0^2$ and $N_0=690988$ (the enclosures are computed over
the rational interval $[t_0,\,t_0+10^{-9}]$, which contains $t_0$);
it reads no stored data, and every acceptance condition below is a
strict inequality between exact rational endpoints, checked in
exact arithmetic.

A single evaluation dominates the whole region $x\ge x_*$,
$y_0\le y\le\sqrt{1-2t_0}$, $t=t_0$ by the following monotonicity
reductions, whose hypotheses are among the twelve domain and sign
conditions certified by the program. Write
$\ell=\log\frac{x_*}{4\pi}=\log\bigl(N_0^2-\tfrac{t_0}{16}\bigr)$,
\[
\delta_1=\frac{t_0}4\log\Bigl(1-\frac{t_0}{16N_0^2}\Bigr)^{-1}
+\frac{t_0}{2x_*^2},
\qquad
\sigma_1=\frac{1+y_0}2+\frac{t_0}2\log N_0-\delta_1
=\frac{1+y_0}2+\frac{t_0}4\ell-\frac{t_0}{2x_*^2}.
\]

(i) \emph{The exponent.} In \eqref{eq:resbound} the positive part
is at most $1$: the program certifies $8/x_*^2<3y_0$, so
$\frac{4y(1+y)}{x^2}\le\frac{8y}{x^2}<3y$ for all $y\ge y_0$, and
$\Rea s_*\ge\frac{1+y}2+\frac{t_0}4\log\frac x{4\pi}
-\frac{t_0}{2x^2}\ge\sigma_1$ throughout the region.

(ii) \emph{The bound \eqref{eq:ec0} for $e_{C,0}$.} Each factor is
monotone in the favorable direction: the prefactor
$(\frac x{4\pi})^{-(1+y)/4}$ and the term
$-\frac{t_0}{16}\log^2\frac x{4\pi}$ decrease in $x$ (certified:
$\ell>2>0$) and in $y$; the factor $3^y+3^{-y}$ increases in $y$;
$N-0.125$ increases with the window; and the final quotient
decreases in $x$ because
$3\,\frac{d}{dx}\bigl|\log\tfrac x{4\pi}+\tfrac{i\pi}2\bigr|
\le\frac3x$ while the numerator exceeds $10.50>3$ (certified).
Hence \eqref{eq:ec0} evaluated at $x=x_*$ and $N=N_0$, with $y=y_0$
in the decreasing factors and $y=\sqrt{1-2t_0}$ in $3^y+3^{-y}$,
bounds $e_{C,0}$ on the whole region; this evaluation is the
program's quantity \texttt{error\_c0}.

(iii) \emph{The bound (23) for $e_A+e_B$.} Since
$4\pi n^2\le4\pi N^2\le x\,(1-\frac{t_0}{16N^2})^{-1}$ for every
$n\le N$, the bounds \eqref{eq:gammabound} and \eqref{eq:kappabound}
give, uniformly on the region,
\[
1+|\gamma|N^{|\kappa|}n^y\ \le\
1+e^{0.02}\Bigl(1-\frac{t_0}{16N_0^2}\Bigr)^{-1/2}
\exp\Bigl(\frac{t_0\log N_0}{2(x_*-6)}\Bigr)
\]
(the quantity $\log N/(x_N-6)$ decreases in $N$), while
$\log^2\frac{x}{4\pi n^2}\le\log^2\frac x{4\pi}$ for $n\le N$ and
the monotonicity of $\log^2(\frac x{4\pi})/(x-6.66)$ in $x$
(certified: $\ell>2$) bound the last factor of (23) by
$\exp\bigl(\frac{(t_0^2/16)\ell^2+0.626}{x_*-6.66}\bigr)-1$. For
the middle factor, $\sum_{n=1}^{N}b_{t_0}(n)\,n^{-\Rea s_*}\le
1+\sum_{n=2}^{m_0}b_{t_0}(n)\,n^{-\sigma_1}+T_N$ with $m_0=2000$,
where $T_N$ is the sum over $m_0<n\le N$. The summand
$\exp(\frac{t_0}4\log^2n-\sigma_N\log n)$ decreases in $n$ on
$n\le N$ (certified: $\frac{1+y_0}2-\delta_1>0$), so
$T_N\le\int_{\log m_0}^{\log N}e^{\psi(L)}\,dL$ with
$\psi(L)=(1-\sigma_N)L+\frac{t_0}4L^2$; since $\psi$ is convex,
$T_N\le(\log N-\log m_0)\,
\max\bigl(e^{\psi(\log m_0)},e^{\psi(\log N)}\bigr)$. Finally, the
two certified rate inequalities
\[
\frac{t_0}2\log m_0\ >\ \Bigl(\log\frac{N_0}{m_0}\Bigr)^{-1},
\qquad
\frac{t_0}2\log N_0-\frac{1-y_0}2-\delta_1\ >\
\Bigl(\log\frac{N_0}{m_0}\Bigr)^{-1},
\]
imply that each of the two endpoint products
$(\log N-\log m_0)e^{\psi(\log m_0)}$ and
$(\log N-\log m_0)e^{\psi(\log N)}$ is nonincreasing in $N$ for
$N\ge N_0$ (using that $\sigma_N$ grows in $\log N$ at rate at
least $\frac{t_0}2$, certified via $N_0^2/(N_0^2-\frac{t_0}{16})>1$,
and that $\delta_N\le\delta_1$), so their values at $N=N_0$
dominate every window. The product of the three bounds is the
program's quantity \texttt{error\_ab}.

The program certifies the twelve domain and sign conditions
required above, computes the enclosures of
\texttt{error\_ab}, \texttt{error\_c0} and of their interval sum,
and establishes the three displayed inequalities by exact rational
comparison of the outward-rounded upper endpoints against the
displayed constants. The run completes in under a second; a
human-readable report of all intermediate enclosures is written to
\texttt{code/prop410/runs/prop410\_report.txt}.
\end{proof}

We can now state the conclusion of the finite range.

\begin{proposition}[Nonvanishing at time $t_0$, finite range]
\label{prop:finitelane}
$H_{t_0}(x+iy)\neq0$ for all $x_*\le x<x_{3840001}$ and
$y_0\le y\le\sqrt{1-2t_0}$.
\end{proposition}

\begin{proof}
Let $x$ be in the stated range and let $N$ be its window index, so
that $690988\le N\le3840000$ by \eqref{eq:xstarwindow}, and
$x\in W_N$. For $y$ in the stated range, Lemma~\ref{lem:native}
(applied with the height $y$, the prime set $\calP(N)$, and the
frozen constants of Lemma~\ref{lem:freeze}), together with
Propositions~\ref{prop:sweep} and~\ref{prop:transfer}, gives
\[
|f_{t_0}(x+iy)|\ \ge\ L_N(y)\ \ge\ L_N(y_0)\ \ge\ T_{\min}.
\]
Combining with Proposition~\ref{prop:finiteerr},
\begin{equation}\label{eq:finitemargin}
|f_{t_0}(x+iy)|-\bigl(e_A+e_B+e_{C,0}\bigr)
\ \ge\ T_{\min}-E_{\max}
\ \ge\ 0.000000557871094787\ >\ 0 .
\end{equation}
The point $(x,y,t_0)$ lies in the region \eqref{eq:region}, so
\cite[Corollary~1.4]{Polymath} (the criterion for non-vanishing
quoted in Section~\ref{sec:thm13}) yields
$H_{t_0}(x+iy)\neq0$.
\end{proof}

% ===============================================================
\section{Hypothesis (ii), infinite tail}\label{sec:tail}
% ===============================================================

This section proves nonvanishing of $H_t$ for \emph{all} windows
$N\ge N_*=3840000$, all
$t\in I_t=\bigl[\frac{129}{800},\frac{161250001}{10^9}\bigr]$ (a
closed interval whose left endpoint is exactly $t_0$), and all
$y\in[y_0,\sqrt{1-2t}]$, by a single interval computation at the
cutoff $N_*$ combined with monotonicity arguments in $N$ and $y$.
Since $t_0\in I_t$, this in particular covers the case $t=t_0$
needed for hypothesis (ii); the last window $W_{3840000}$ of the
finite range overlaps the tail's starting abscissa $x_{3840000}$, so
the two ranges jointly cover all $x\ge x_*$.

\subsection{Setup: mollifier and frozen exponents}

Fix the sparse mollifier support
$\calS=\{1,2,3,4,5,6,7,10,11,13,14\}$ and the exact rational vector
\[
\begin{array}{c|ccccccccccc}
d&1&2&3&4&5&6&7&10&11&13&14\\ \hline
\lambda_d&
1&-\frac{1021}{1000}&-\frac{1054}{1000}&-\frac{9}{200}&
-\frac{1119}{1000}&\frac{1001}{1000}&-\frac{1043}{1000}&
\frac{128}{125}&-\frac{161}{200}&-\frac{447}{500}&\frac{456373}{500000}
\end{array}
\]
($\lambda_1=1$; $\sum_d|\lambda_d|=9.918746$ exactly), and set
$M_\lambda(z)=\sum_{d\in\calS}\lambda_d d^{-z}$. (Unlike the Euler-product
mollifier of Section~\ref{sec:native}, this vector was chosen by numerical optimization;
its only structural requirement is $\lambda_1=1$ and real entries.)
Fix also the head lengths $M=153814$ and $M_{\rm err}=3000$.

For $x\ge X_N(t):=4\pi\bigl(N^2-\frac t{16}\bigr)$ with window index
$N\ge N_*$, define frozen exponents (uniform over $I_t$, via rigorous
upper bounds $\widehat\delta\ge\delta(N_*,t)$,
$\widehat k\ge k(N_*,t)$ for the quantities
$\delta(N,t)=\frac t4\bigl[-\log\bigl(1-\frac t{16N^2}\bigr)\bigr]
+\frac{t}{2X_N(t)^2}$ and $k(N,t)=\frac{t}{2(X_N(t)-6)}$, both
decreasing in $N$):
\[
\sigma_1=\frac{1+y}{2}+\frac t2\log N-\widehat\delta,
\qquad
\sigma_2=\frac{1-y}{2}+\frac t2\log N-\widehat\delta-\widehat k,
\qquad
G=e^{0.02y}\Bigl(N^2-\frac t{16}\Bigr)^{-y/2}.
\]
\begin{lemma}[Frozen exponents]\label{lem:frozen}
Let $N\ge N_*$, $t\in I_t$, $0<y\le1$, and let
$x\in[X_N(t),X_{N+1}(t))$, so that $x$ has window index $N$. Then
\[
|\gamma|\le G,\qquad
\Rea s_*\ge\sigma_1,\qquad
|\kappa|\le\widehat k,
\]
and consequently each term of the second sum in \eqref{eq:ft}
satisfies
\[
\bigl|n^y\,b_t(n)\,n^{-\overline{s_*}-\kappa}\bigr|
\ \le\ b_t(n)\,n^{-\sigma_2}
\qquad(1\le n\le N).
\]
\end{lemma}

\begin{proof}
Since $x\ge X_N(t)$ we have $x/4\pi\ge N^2-t/16$, and the three
factors $(x/4\pi)^{-y/2}$, $x^{-2}$, $(x-6)^{-1}$ appearing in
\eqref{eq:gammabound}--\eqref{eq:kappabound} are decreasing in $x$,
hence maximal at the window's
left edge. For $\gamma$ this gives
$|\gamma|\le e^{0.02y}(N^2-t/16)^{-y/2}=G$ directly. For $\kappa$,
using $y\le1$ and that $X_N(t)$ increases with $N$,
\[
|\kappa|\ \le\ \frac{ty}{2(x-6)}\ \le\ \frac{t}{2(X_N(t)-6)}
\ =\ k(N,t)\ \le\ k(N_*,t)\ \le\ \widehat k .
\]
For $\Rea s_*$: by the exact identity
\[
\frac t4\log\Bigl(N^2-\frac t{16}\Bigr)
=\frac t2\log N
-\frac t4\Bigl[-\log\Bigl(1-\frac t{16N^2}\Bigr)\Bigr],
\]
the logarithmic term of \eqref{eq:resbound} is at least
$\frac t2\log N$ minus the first summand of $\delta(N,t)$. The
subtracted positive-part term of \eqref{eq:resbound} is at most
$t/(2x^2)\le t/(2X_N(t)^2)$, the second summand of $\delta(N,t)$:
indeed, on the domain $0\le y\le1$, $x\ge200$ one has
$1-3y+4y(1+y)/x^2\le1$ (for $y>0$ this is $4(1+y)\le3x^2$, and at
$y=0$ the left side equals $1$), so the positive part is at most
$1$. Both summands of $\delta(N,t)$ decrease in $N$ (in the first,
$t/16N^2$ decreases and $u\mapsto-\log(1-u)$ is increasing; in the
second, $X_N(t)$ increases), so
\[
\Rea s_*\ \ge\ \frac{1+y}2+\frac t2\log N-\delta(N,t)
\ \ge\ \frac{1+y}2+\frac t2\log N-\widehat\delta\ =\ \sigma_1 .
\]
Finally, for $1\le n\le N$,
\[
\bigl|n^y b_t(n)n^{-\overline{s_*}-\kappa}\bigr|
= b_t(n)\,n^{\,y-\Rea s_*-\Rea\kappa}
\le b_t(n)\,n^{\,y-\sigma_1+\widehat k}
= b_t(n)\,n^{-\sigma_2},
\]
since $-\Rea\kappa\le|\kappa|\le\widehat k$ and
$\sigma_2=\sigma_1-y-\widehat k$.
\end{proof}

\subsection{The endpoint-cap lemma}

The analytic engine for tail sums is:

\begin{lemma}[Endpoint cap]\label{lem:cap}
For $1\le a<c$ set
$E_{t,\sigma}(u)=\exp\bigl((1-\sigma)\log u+\frac t4\log^2u\bigr)$ and
\[
\CapT_t(a,c;\sigma)=\max\bigl(E_{t,\sigma}(a),E_{t,\sigma}(c)\bigr)\log\frac ca .
\]
If $\sigma>\frac t2\log c$, then for all integers $a\le J<K\le c$,
\[
\sum_{J<n\le K}b_t(n)\,n^{-\sigma}\ \le\ \CapT_t(a,c;\sigma).
\]
\end{lemma}

\begin{proof}
The logarithmic derivative of $b_t(u)u^{-\sigma}$ in $\log u$ is
$\frac t2\log u-\sigma<0$ on $[a,c]$, so the summand decreases and the
sum is at most $\int_a^c b_t(u)u^{-\sigma}du$. Substituting
$v=\log u$ makes the integrand $\exp\bigl((1-\sigma)v+\frac t4v^2\bigr)$,
which is log-convex in $v$, hence maximized at an endpoint of
$[\log a,\log c]$; multiplying by the length $\log(c/a)$ of the
$v$-interval gives the bound.
\end{proof}

\subsection{Exact convolution and the contraction $D$}

Define, for $N\ge N_*$ and $(x,y,t)$ with $x$ in the window of index
$N$,
\[
P=\sum_{m=2}^{M}|c_m(t)|m^{-\sigma_1},
\qquad
c_m(t)=\sum_{\substack{d\in\calS\\ d\mid m}}\lambda_d\,b_t(m/d),
\qquad
M_{\max}=\sum_{d\in\calS}|\lambda_d|\,d^{-\sigma_1},
\]
\[
\mathrm{TR}=\sum_{d\in\calS}|\lambda_d|\,d^{-\sigma_1}
\CapT_t\bigl(\lfloor M/d\rfloor,N;\sigma_1\bigr),
\qquad
\mathrm{AB}=G\Bigl(\sum_{k\le M}b_t(k)k^{-\sigma_2}
+\CapT_t(M,N;\sigma_2)\Bigr),
\]
together with the nonnegative ``overshoot'' term
\[
\mathrm{OV}=\sum_{\substack{d\in\calS\\ d>1}}
|\lambda_d|\,d^{-\sigma_1}\,
\CapT_t\Bigl(\frac{N}{d+1},\,N;\sigma_1\Bigr)
\]
(included by Gomila's implementation; see the remark after the
proof), and set
\begin{equation}\label{eq:D}
D\ =\ D(N,t,y)\ :=\ P+\mathrm{TR}+\mathrm{OV}+M_{\max}\,\mathrm{AB}.
\end{equation}

\begin{lemma}[Mollified contraction bound]\label{lem:contraction}
Assume the cap-validity conditions
\begin{equation}\label{eq:capvalid}
\widehat\delta\ <\ \frac{1+y}2,
\qquad
\widehat\delta+\widehat k\ <\ \frac{1-y}2
\end{equation}
(equivalently $\sigma_1>\frac t2\log N$ and
$\sigma_2>\frac t2\log N$: the differences
$\sigma_i-\frac t2\log N$ are independent of $N$). Then for every
$N\ge N_*$, $t\in I_t$, $0<y\le1$ and $x$ in the window of index
$N$,
\[
\bigl|M_\lambda(s_*)f_t(x+iy)-1\bigr|\ \le\ D .
\]
Consequently, if $D<1$ then
\begin{equation}\label{eq:flowlower}
|f_t(x+iy)|\ \ge\ \frac{1-D}{M_{\max}} .
\end{equation}
\end{lemma}

\begin{proof}
Write the two sums in \eqref{eq:ft} as $F_B$ and $F_A$
($f_t=F_B+\gamma F_A$). Expanding
$M_\lambda(s_*)F_B=\sum_{d\in\calS}\sum_{k\le N}
\lambda_d\,b_t(k)\,(dk)^{-s_*}$ and partitioning the index pairs
$(d,k)$ by $dk\le M$ versus $dk>M$ --- a disjoint and exhaustive
partition --- gives the \emph{exact} identity
\begin{equation}\label{eq:tailconv}
M_\lambda(s_*)F_B-1
=\sum_{m=2}^{M}c_m(t)\,m^{-s_*}
+\sum_{d\in\calS}\lambda_d\,d^{-s_*}
\sum_{\lfloor M/d\rfloor<k\le N}b_t(k)\,k^{-s_*} .
\end{equation}
Here the first term groups the pairs with $dk\le M$ by $m=dk$: since
$m\le M<N_*\le N$, the constraint $k=m/d\le N$ is automatic, so the
coefficient of $m^{-s_*}$ is exactly $c_m(t)$, and the $m=1$ term
$c_1=\lambda_1b_t(1)=1$ has been moved to the left side. In the
second term, for fixed $d$ the condition $dk>M$ for integer $k$ is
$k>\lfloor M/d\rfloor$. By $\Rea s_*\ge\sigma_1$
(Lemma~\ref{lem:frozen}), the first term is bounded by $P$; and
each inner sum of the second term is bounded, by
Lemma~\ref{lem:cap} with $(a,c)=(\lfloor M/d\rfloor,N)$, by
$\CapT_t(\lfloor M/d\rfloor,N;\sigma_1)$, so the second term is
bounded by $\mathrm{TR}$; the application of Lemma~\ref{lem:cap} is
legitimate because its hypothesis $\sigma_1>\frac t2\log N$ is the
first condition of \eqref{eq:capvalid}. On the $A$-side,
Lemma~\ref{lem:frozen} bounds each term of $|\gamma F_A|$ by
$G\,b_t(k)k^{-\sigma_2}$; splitting at $k=M$ and applying
Lemma~\ref{lem:cap} with $(a,c)=(M,N)$ at exponent $\sigma_2$
(legitimate by the second condition of \eqref{eq:capvalid}) gives
$|\gamma F_A|\le\mathrm{AB}$, while
$|M_\lambda(s_*)|\le\sum_d|\lambda_d|d^{-\sigma_1}=M_{\max}$;
hence $|M_\lambda(s_*)\gamma F_A|\le M_{\max}\mathrm{AB}$. Adding
the nonnegative $\mathrm{OV}$ only weakens the bound, and the first
claim follows. For the second claim: if $D<1$ then
$|M_\lambda(s_*)f_t|\ge1-D>0$; in particular
$M_\lambda(s_*)\neq0$, and dividing by
$|M_\lambda(s_*)|\le M_{\max}$ gives
$|f_t|\ge(1-D)/M_{\max}$.
\end{proof}

\begin{remark}
The term $\mathrm{OV}$ plays no role in the mathematics --- the
identity \eqref{eq:tailconv} is fully covered by $P+\mathrm{TR}$ ---
but Gomila's program includes it in $D$, so it is included in
\eqref{eq:D} in order that the certified numerical bound on $D$
(Proposition~\ref{prop:tailcert} below) majorizes the same quantity
the program computes.
\end{remark}

\subsection{Reduction to the cutoff}\label{sec:tailmono}

The interval computation of Section~\ref{sec:tailcert} evaluates
$D$, and the error majorants, only at the cutoff $N=N_*$, over the
closed intervals
\[
I_{\rm box}=\Bigl[\tfrac{1872719}{10^7},\tfrac{23409}{125000}\Bigr],
\qquad
I_{\rm ext}=\Bigl[\tfrac{1872719}{10^7},\tfrac{8231039}{10^7}\Bigr],
\]
where $I_{\rm box}$ contains $y_0$ and the top of $I_{\rm ext}$ is
at least $\sqrt{1-2t}$ for every $t\in I_t$. The reduction from
this single evaluation to all $N\ge N_*$ and the full height range
is carried out in Lemmas~\ref{lem:capdecay}--\ref{lem:errmaj}
below. Throughout, the \emph{cutoff conditions} are the following
inequalities, all evaluated at $N=N_*$, uniformly for $t\in I_t$:
\begin{equation}\label{eq:gates}
\frac t2\log a\,(\log N_*-\log a)>1
\quad\text{and}\quad
(\sigma-1)(\log N_*-\log a)>1
\end{equation}
for each cap $\CapT_t(a,\,\cdot\,;\sigma)$ appearing in $D$ and in
the error bounds (the left endpoints $a=\lfloor M/d\rfloor$ for
$d\in\calS$, $a=M$, and $a=M_{\rm err}$, with $\sigma=\sigma_1$ or
$\sigma_2$ as appropriate), together with the $A$-side slope
condition
\begin{equation}\label{eq:ym}
0.02-\tfrac12\log\bigl(N_*^2-\tfrac t{16}\bigr)+\tfrac12\log N_*\ <\ 0 .
\end{equation}

\begin{lemma}[Decay of fixed-endpoint caps]\label{lem:capdecay}
Let $t>0$, let $1<a<N_*$, let $\sigma(q)=\sigma_0+\frac t2q$ with
$\sigma_0\in\R$, and write $q_*=\log N_*$. Assume the two
conditions \eqref{eq:gates} for this pair, that is,
$\frac t2\log a\,(q_*-\log a)>1$ and
$(\sigma(q_*)-1)(q_*-\log a)>1$. Then both conditions persist for
all $q\ge q_*$, one has $\sigma(q)>1$ for all $q\ge q_*$, and
\[
q\ \longmapsto\ \CapT_t\bigl(a,e^q;\sigma(q)\bigr)
\]
is positive and decreasing on $[q_*,\infty)$.
\end{lemma}

\begin{proof}
By definition
$\CapT_t(a,e^q;\sigma(q))=\max\bigl(F_1(q),F_2(q)\bigr)$, where
\[
F_1(q)=E_{t,\sigma(q)}(a)\,(q-\log a),
\qquad
F_2(q)=E_{t,\sigma(q)}(e^q)\,(q-\log a).
\]
Since $\log E_{t,\sigma}(u)=(1-\sigma)\log u+\frac t4\log^2u$ and
$\sigma'(q)=\frac t2$,
\[
\frac{d}{dq}\log F_1(q)
=-\frac t2\log a+\frac1{q-\log a},
\]
which is negative exactly when
$\frac t2\log a\,(q-\log a)>1$; both factors on the left of this
condition are positive and nondecreasing in $q$, so the condition,
true at $q_*$, persists for all $q\ge q_*$. Similarly,
\[
\frac{d}{dq}\log F_2(q)
=\Bigl[(1-\sigma(q))-\frac t2\,q\Bigr]+\frac t2\,q
+\frac1{q-\log a}
=1-\sigma(q)+\frac1{q-\log a},
\]
which is negative exactly when $(\sigma(q)-1)(q-\log a)>1$. At
$q_*$ this condition forces $\sigma(q_*)>1$; then both factors
$\sigma(q)-1$ and $q-\log a$ are positive and increasing in $q$,
so the condition persists, and $\sigma(q)>1$ throughout. Hence
$F_1$ and $F_2$ are positive and decreasing on $[q_*,\infty)$, and
so is their maximum.
\end{proof}

\begin{lemma}[Componentwise decay in $N$]\label{lem:compdecay}
Fix $t\in I_t$ and $y\in(0,1]$, and assume the cutoff conditions
\eqref{eq:gates}. Then each of the following is nonincreasing in
$N$ on $[N_*,\infty)$ (throughout, $q=\log N$ and the frozen
exponents $\sigma_i=\sigma_i(q)$ are increasing in $q$ with slope
$\frac t2$):
\begin{enumerate}
\item[\textup{(i)}] $m^{-\sigma_i(q)}$ for every fixed integer
$m\ge2$ and $i\in\{1,2\}$; hence the head $P$, the mollifier
bound $M_{\max}$, and the $A$-side head
$\sum_{k\le M}b_t(k)k^{-\sigma_2}$ (whose $k=1$ term is constant)
are nonincreasing;
\item[\textup{(ii)}] $G=e^{0.02y}\bigl(e^{2q}-\frac t{16}\bigr)^{-y/2}$;
\item[\textup{(iii)}] the moving-endpoint term $\mathrm{OV}$.
\end{enumerate}
\end{lemma}

\begin{proof}
(i) is immediate: $\sigma_i(q)$ is increasing and $\log m>0$,
while the coefficients $|c_m(t)|$, $|\lambda_d|$ and $b_t(k)$ do
not depend on $N$. For (ii),
\[
\frac{d}{dq}\log G
=-\frac y2\cdot\frac{2e^{2q}}{e^{2q}-t/16}
=-\,y\,\frac{N^2}{N^2-t/16}\ <\ 0 .
\]
For (iii), each summand of $\mathrm{OV}$ is
$|\lambda_d|\,d^{-\sigma_1(q)}\,
\CapT_t(e^q/(d+1),e^q;\sigma_1(q))$, a product of the
nonincreasing factor $d^{-\sigma_1(q)}$ ($d\ge2$) with a cap whose
endpoints both scale with $N$ and whose length $\log(d+1)$ is
constant. For $r\in\{1,d+1\}$,
\[
\frac{d}{dq}\Bigl[(1-\sigma(q))(q-\log r)
+\frac t4(q-\log r)^2\Bigr]
=(1-\sigma(q))-\frac t2(q-\log r)+\frac t2(q-\log r)
=1-\sigma(q)\ <\ 0,
\]
using $\sigma_1(q)>1$ from Lemma~\ref{lem:capdecay}. So both cap
candidates are positive and decreasing, hence so are their maximum
and the summand.
\end{proof}

\begin{lemma}[Decay in $y$ at the cutoff]\label{lem:ydecay}
Fix $t\in I_t$ and assume the $A$-side slope condition
\eqref{eq:ym}. Then $y\mapsto D(N_*,t,y)$ and
$y\mapsto M_{\max}(N_*,t,y)$ are nonincreasing on $(0,1]$.
\end{lemma}

\begin{proof}
Recall $\partial\sigma_1/\partial y=\frac12$ and
$\partial\sigma_2/\partial y=-\frac12$. On the $B$-side, every
$y$-dependent factor has the form $u^{-\sigma_1(y)}$ or
$E_{t,\sigma_1(y)}(u)$ with $u>1$ (the head terms $m\ge2$, the
mollifier terms $d\ge2$, and the cap ingredients at
$u=\lfloor M/d\rfloor$, $u=N_*$ or $u=N_*/r$, all exceeding $1$),
and satisfies $\partial_y\log(\cdot)=-\frac12\log u<0$; the
remaining factors ($|c_m(t)|$, $|\lambda_d|$, $b_t$, cap lengths)
are independent of $y$, and maxima of nonincreasing functions are
nonincreasing. Hence $P$, $\mathrm{TR}$, $\mathrm{OV}$ and
$M_{\max}$ are nonincreasing in $y$. On the $A$-side, $\log G$ is
linear in $y$ with slope
$0.02-\frac12\log\bigl(N_*^2-\frac t{16}\bigr)$, so each head
term and each cap ingredient of $\mathrm{AB}$, with $u\le N_*$,
satisfies
\[
\partial_y\log\Bigl(G\,b_t(u)\,u^{-\sigma_2(y)}\Bigr)
\ \text{resp.}\ \
\partial_y\log\Bigl(G\,E_{t,\sigma_2(y)}(u)\Bigr)
\ =\ 0.02-\tfrac12\log\Bigl(N_*^2-\tfrac t{16}\Bigr)
+\tfrac12\log u,
\]
which is at most the left side of \eqref{eq:ym}, hence negative.
So $\mathrm{AB}$ is nonincreasing in $y$, and therefore so are
$M_{\max}\mathrm{AB}$ and $D$.
\end{proof}

\begin{lemma}[Reduction to the cutoff]\label{lem:reduction}
Assume the cutoff conditions \eqref{eq:gates}--\eqref{eq:ym} hold.
Then for every $N\ge N_*$, every $t\in I_t$ and every
$y\in[y_0,\sqrt{1-2t}]$,
\[
D(N,t,y)\ \le\ \sup_{y'\in I_{\rm box}}D(N_*,t,y'),
\qquad
M_{\max}(N,t,y)\ \le\ \sup_{y'\in I_{\rm box}}M_{\max}(N_*,t,y').
\]
\end{lemma}

\begin{proof}
Fix $t$ and $y$. Every constituent of
$D=P+\mathrm{TR}+\mathrm{OV}+M_{\max}\mathrm{AB}$ is a sum of
products of nonnegative factors, each nonincreasing in $N$ on
$[N_*,\infty)$: the heads, $M_{\max}$ and $G$ by
Lemma~\ref{lem:compdecay}(i)--(ii); the fixed-endpoint caps in
$\mathrm{TR}$ and $\mathrm{AB}$ by Lemma~\ref{lem:capdecay},
applied with $a=\lfloor M/d\rfloor$ (at $\sigma_1$) and $a=M$ (at
$\sigma_2$), whose conditions \eqref{eq:gates} hold by hypothesis;
and $\mathrm{OV}$ by Lemma~\ref{lem:compdecay}(iii). Sums and
products of nonnegative nonincreasing functions are nonincreasing,
so $D(N,t,y)\le D(N_*,t,y)$ and
$M_{\max}(N,t,y)\le M_{\max}(N_*,t,y)$. If $y\in I_{\rm box}$
the claims follow at once. Otherwise $y>\max I_{\rm box}$, since
$y\ge y_0\ge\min I_{\rm box}$; as $y\le\sqrt{1-2t}<1$,
Lemma~\ref{lem:ydecay} applies on $[\max I_{\rm box},y]$ and
gives $D(N_*,t,y)\le D(N_*,t,\max I_{\rm box})
\le\sup_{y'\in I_{\rm box}}D(N_*,t,y')$, and the same for
$M_{\max}$.
\end{proof}

\subsection{Tail error terms}\label{sec:tailerrors}

The error bounds use the exact definitions (71)--(72) of
\cite{Polymath} rather than the coarser display (23): (23) applies
the worst-case factor $N^{|\kappa|}$ to every term, whereas (71)
carries the per-term factor $n^{-\Re\kappa}$, which
Lemma~\ref{lem:frozen} bounds by $n^{\widehat k}$. Define, for
$N\ge N_*$, $t\in I_t$ and $0<y\le1$,
\[
L_N(t)=\log\Bigl(N^2-\frac t{16}\Bigr),
\qquad
\Delta(N,t)=\frac{\frac{t^2}{16}L_N(t)^2+0.626}{X_N(t)-6.66},
\]
and the two majorants
\begin{align*}
\mathrm{E}_{AB}(N,t,y)&=
\bigl(e^{\Delta(N,t)}-1\bigr)
\bigl(S_1(N,t,y)+G\,S_2(N,t,y)\bigr),
\\
S_i(N,t,y)&=\sum_{n=1}^{M_{\rm err}}b_t(n)\,n^{-\sigma_i}
+\CapT_t\bigl(M_{\rm err},N;\sigma_i\bigr)
\qquad(i=1,2),
\\
\mathrm{E}_{C}(N,t,y)&=
\Bigl(N^2-\frac t{16}\Bigr)^{-(1+y)/4}
\exp\biggl(
-\frac t{16}L_N(t)^2
+\frac{1.24\,(3^y+3^{-y})}{N-0.125}
+\frac{3\sqrt{L_N(t)^2+\pi^2/4}+10.50}{X_N(t)-12}
\biggr).
\end{align*}

\begin{lemma}[Error majorants]\label{lem:errmaj}
Assume the cap-validity conditions \eqref{eq:capvalid} and the
conditions \eqref{eq:gates} for the caps with left endpoint
$a=M_{\rm err}$. Then for every $N\ge N_*$, $t\in I_t$,
$y\in[y_0,\sqrt{1-2t}]$ and $x$ in the window of index $N$,
\[
e_A+e_B\ \le\ \mathrm{E}_{AB}(N,t,y)\ \le\ \mathrm{E}_{AB}(N_*,t,y),
\qquad
e_{C,0}\ \le\ \mathrm{E}_{C}(N,t,y)\ \le\ \mathrm{E}_{C}(N_*,t,y).
\]
\end{lemma}

\begin{proof}
\emph{The per-term factor.} The quantity $\eps_{t,n}$ of
\cite[(44)]{Polymath} is bounded, by the display in the proof of
\cite[Proposition~6.6(iv),(v)]{Polymath},
\[
\eps_{t,n}\Bigl(\frac{1\pm y+ix}2\Bigr)
\ \le\
\exp\biggl(
\frac{\frac{t^2}{32}\log^2\frac{x}{4\pi n^2}+0.313}
{\frac x2-3.33}
\biggr)-1
\ =\
\exp\biggl(
\frac{\frac{t^2}{16}\log^2\frac{x}{4\pi n^2}+0.626}{x-6.66}
\biggr)-1 ,
\]
the imaginary part of $\frac{1\pm y+ix}2$ being $\frac x2$. For
$1\le n\le N$ and $x$ in the window of index $N$ one has
$\frac{x}{4\pi n^2}\le\frac{x}{4\pi}$ and
$\frac{x}{4\pi n^2}\ge\frac{X_N(t)}{4\pi N^2}
=1-\frac t{16N^2}$, so
\[
\Bigl|\log\frac{x}{4\pi n^2}\Bigr|
\ \le\ \max\Bigl(\log\frac x{4\pi},
\,-\log\Bigl(1-\frac t{16N^2}\Bigr)\Bigr)
\ =\ \log\frac x{4\pi},
\]
the last step because
$\frac x{4\pi}\bigl(1-\frac t{16N^2}\bigr)
\ge\frac{(N^2-t/16)^2}{N^2}\ge1$ (indeed
$N^2-\frac t{16}\ge N$ as $N\ge N_*$). Hence the per-term
exponent is at most
$\varphi(x)
=\bigl(\frac{t^2}{16}\log^2\frac x{4\pi}+0.626\bigr)/(x-6.66)$.
Writing $L=\log\frac x{4\pi}$, which exceeds $2$ for $x\ge200$,
\[
(x-6.66)^2\,\varphi'(x)
=\frac{t^2}{8}\,L\,\frac{x-6.66}{x}
-\Bigl(\frac{t^2}{16}L^2+0.626\Bigr)
\ <\ \frac{t^2}{8}L-\frac{t^2}{16}L^2-0.626
=-\frac{t^2}{16}L(L-2)-0.626\ <\ 0,
\]
so $\varphi$ is decreasing, and on the window
$\varphi(x)\le\varphi(X_N(t))=\Delta(N,t)$: the per-term factor
is at most $e^{\Delta(N,t)}-1$.

\emph{Assembly of $\mathrm{E}_{AB}$.} By (71)--(72) of
\cite{Polymath} and Lemma~\ref{lem:frozen}
($\Rea s_*\ge\sigma_1$, $|\gamma|\le G$, and
$n^{\,y-\Rea s_*-\Rea\kappa}\le n^{-\sigma_2}$),
\[
e_B\le\bigl(e^{\Delta(N,t)}-1\bigr)
\sum_{n\le N}b_t(n)\,n^{-\sigma_1},
\qquad
e_A\le G\bigl(e^{\Delta(N,t)}-1\bigr)
\sum_{n\le N}b_t(n)\,n^{-\sigma_2}.
\]
Splitting each sum at $M_{\rm err}$ and bounding the range
$M_{\rm err}<n\le N$ by Lemma~\ref{lem:cap} (legitimate by
\eqref{eq:capvalid}) gives
$e_A+e_B\le\mathrm{E}_{AB}(N,t,y)$.

\emph{Bound for $e_{C,0}$.} By the corollary \eqref{eq:ec0},
$e_{C,0}$ is bounded by the product of
$(\frac x{4\pi})^{-(1+y)/4}$,
$\exp(-\frac t{16}\log^2\frac x{4\pi})$,
$\exp\bigl(1.24(3^y+3^{-y})/(N-0.125)\bigr)$, and
$\exp(\psi(x))$ with
$\psi(x)=\bigl(3\sqrt{L^2+\pi^2/4}+10.50\bigr)/(x-12)$. The
first two factors are decreasing in $x$, the third does not depend
on $x$, and $\psi$ is decreasing: with
$Q(L)=3\sqrt{L^2+\pi^2/4}+10.50$ one has
$\frac{d}{dx}Q=\frac{3L}{\sqrt{L^2+\pi^2/4}}\cdot\frac1x
<\frac3x$, whence
$(x-12)^2\,\psi'(x)<3\,\frac{x-12}x-Q<3-10.50<0$. Evaluating
all four factors at the left edge $x=X_N(t)$, where
$\log\frac x{4\pi}=L_N(t)$, gives
$e_{C,0}\le\mathrm{E}_{C}(N,t,y)$.

\emph{Decay in $N$.} In $\mathrm{E}_{AB}$: the factor
$e^{\Delta(N,t)}-1$ is nonincreasing because
$\Delta(N,t)=\varphi(X_N(t))$ with $\varphi$ decreasing and
$X_N(t)$ increasing in $N$; the heads of $S_1,S_2$ and the factor
$G$ are nonincreasing by Lemma~\ref{lem:compdecay}; and the caps
$\CapT_t(M_{\rm err},N;\sigma_i)$ are nonincreasing by
Lemma~\ref{lem:capdecay} with $a=M_{\rm err}$, whose conditions
\eqref{eq:gates} hold by hypothesis. In $\mathrm{E}_{C}$: the
factors $(N^2-t/16)^{-(1+y)/4}$, $e^{-\frac t{16}L_N(t)^2}$ and
$e^{1.24(3^y+3^{-y})/(N-0.125)}$ are visibly nonincreasing in $N$,
and $e^{\psi(X_N(t))}$ is nonincreasing because $\psi$ is
decreasing and $X_N(t)$ increases. Products of nonnegative
nonincreasing factors are nonincreasing.
\end{proof}

\subsection{The certified cutoff inequalities and the tail theorem}
\label{sec:tailcert}

The computational input of the tail argument is the following.

\begin{proposition}[Certified cutoff inequalities]
\label{prop:tailcert}
The cap-validity conditions \eqref{eq:capvalid} and the cutoff
conditions \eqref{eq:gates}--\eqref{eq:ym} hold, uniformly for
$t\in I_t$ and $y\in[y_0,\sqrt{1-2t}]$, and
\[
\sup_{t\in I_t,\ y\in I_{\rm box}}D(N_*,t,y)<0.999721,
\qquad
M_{\max}<1.608290,
\]
\[
\sup_{t\in I_t,\ y\in I_{\rm ext}}
\bigl(\mathrm{E}_{AB}+\mathrm{E}_{C}\bigr)(N_*,t,y)
<1.1672\times10^{-8},
\]
with $\mathrm{E}_{AB},\mathrm{E}_{C}$ the error majorants of
Lemma~\ref{lem:errmaj}. In particular
\[
\frac{1-D}{M_{\max}}-(e_A+e_B+e_{C,0})\ >\ 0.00017352\ >\ 0 .
\]
\end{proposition}

\begin{proof}
The proof is computer-assisted; the verification program
\texttt{prop510\_proof.c}, included with this manuscript as
supplementary material (directory \texttt{code/prop510}; derived
from Gomila's source as documented in its header and reviewed line
by line, check by check, as part of this audit), certifies every
condition and every displayed constant of the statement. It reads
no stored data; every decisive comparison is a strict ball
inequality, which fails closed when enclosures overlap; and it was
run at both $256$ and $512$ bits of precision with identical
verdicts.

\emph{How the suprema are captured.} The intervals $I_t$,
$I_{\rm box}$ and $I_{\rm ext}$ enter the computation as single
balls whose endpoints are the exact rational interval endpoints;
consequently every derived quantity is an enclosure valid
simultaneously for all $(t,y)$ in the corresponding hull, and its
directed upper (resp.\ lower) endpoint bounds the supremum (resp.\
infimum) over that hull. The frozen quantities $\widehat\delta$ and
$\widehat k$ are extracted as point \emph{upper} bounds over $I_t$,
and the exponents $\sigma_1,\sigma_2$ as point \emph{lower} bounds
over the hulls; this is conservative throughout because every
occurrence of $\sigma_i$ is through $u^{-\sigma_i}$ or
$E_{t,\sigma_i}(u)$ with $u\ge1$, both nonincreasing in $\sigma_i$.

\emph{What is certified.} (1) The cap-validity conditions
\eqref{eq:capvalid}, in the equivalent form
$\sigma_i>\frac t2\log N_*$, on both hulls (labels \texttt{SC1},
\texttt{SC2} in the program's output). (2) The cutoff conditions
\eqref{eq:gates} for exactly the left endpoints listed there:
$a=\lfloor M/d\rfloor$ for each $d\in\calS$ at $\sigma_1$
(\texttt{GN-TR}), $a=M$ at $\sigma_2$ (\texttt{GN-AB}), and
$a=M_{\rm err}$ at both $\sigma_1$ and $\sigma_2$
(\texttt{GN-error}), each with the $t$-factor evaluated at
$\min I_t$ and the $\sigma$-factor at its lower bound over
$I_{\rm ext}$. (3) The $A$-side slope condition \eqref{eq:ym}
(\texttt{YM}). (4) The hypotheses of
Lemmas~\ref{lem:capdecay}--\ref{lem:errmaj}: $\sigma_1,\sigma_2>1$
on $I_{\rm ext}$, $N_*\bigl(1-\frac t{16N_*^2}\bigr)>1$,
$X_{N_*}(t)>200$, $\log(N_*^2-\frac t{16})>2$, and the sign
conditions making $\Delta(N,t)$ and the $e_{C,0}$ quotient
decreasing. (5) The exact rational identities: the mollifier
normalization $\sum_d|\lambda_d|=9.918746$, $\lambda_1=1$,
$t_0+y_0^2/2=\frac{893927}{5\cdot10^6}$, and the hull containments
$y_0\in I_{\rm box}$,
$\max I_{\rm box}\le\sqrt{1-2t}$ for all $t\in I_t$, and
$\max I_{\rm ext}\ge\sqrt{1-2t_0}$, all verified in exact integer
arithmetic. (6) The quantity $D$ is assembled per \eqref{eq:D}
from an exact Dirichlet convolution $c_m(t)$ computed by direct
double summation over $(d,k)$, and the four displayed constants
are certified as strict ball comparisons against exact rationals:
\[
D<0.999721,\qquad M_{\max}<1.608290,\qquad
\mathrm{E}_{AB}+\mathrm{E}_{C}<1.1672\times10^{-8},
\]
\[
\frac{1-D}{M_{\max}}-(\mathrm{E}_{AB}+\mathrm{E}_{C})
\ >\ 0.00017352
\]
(the last check was added in this review: Gomila's original program
certified only positivity of this margin, together with two-sided
``corridor'' checks on each constituent, which are retained since
they expose missing or duplicated terms). The certified enclosures
at $512$ bits include
$D\le0.9997209094$, $M_{\max}\le1.6082891450$,
$\mathrm{E}_{AB}+\mathrm{E}_{C}\le1.16716035\times10^{-8}$, and
margin $\ge0.0001735209$. Both runs ($256$ and $512$ bits) pass all
$37$ checks; the outputs are included as
\texttt{code/prop510/runs/prop510\_report\_256.txt}
and \texttt{prop510\_report\_512.txt}, and each run completes in
under a second.
\end{proof}

\begin{theorem}[Tail theorem]\label{thm:tail}
For every $t\in I_t$, $y\in[y_0,\sqrt{1-2t}]$ and
$x\ge X_{N_*}(t)$,
\[
|f_t(x+iy)|\ >\ e_A+e_B+e_{C,0},
\]
and consequently $H_t(x+iy)\ne0$ on this region.
\end{theorem}

\begin{proof}
Let $N$ be the window index of $x$, so $N\ge N_*$. By
Lemma~\ref{lem:contraction}, Lemma~\ref{lem:reduction} and
Proposition~\ref{prop:tailcert},
\[
|f_t(x+iy)|\ \ge\ \frac{1-D(N,t,y)}{M_{\max}}
\ \ge\ \frac{1-0.999721}{M_{\max}}
\ >\ 0.0001735
\ >\ 1.1672\times10^{-8}
\ \ge\ e_A+e_B+e_{C,0},
\]
using in the second step both clauses of
Lemma~\ref{lem:reduction} with Proposition~\ref{prop:tailcert}
($D<0.999721$ and $M_{\max}<1.608290$), and in the last step
Lemma~\ref{lem:errmaj} together with
Proposition~\ref{prop:tailcert} and
$[y_0,\sqrt{1-2t}]\subseteq I_{\rm ext}$. The point $(x,y,t)$
lies in the region \eqref{eq:region}, so
\cite[Corollary~1.4]{Polymath} gives $H_t(x+iy)\neq0$.
\end{proof}

% ===============================================================
\section{Hypothesis (iii): the barrier certificate}\label{sec:barrier}
% ===============================================================

\subsection{Statement, and reduction to a closed rectangle}

Throughout this section let
\[
R=[X,\,X+1]+i\Bigl[\tfrac{1809}{10000},\,1\Bigr].
\]
The goal of the section is the following theorem, whose proof is
assembled in Section~\ref{sec:prismcert}.

\begin{theorem}[Barrier theorem]\label{thm:barrier}
$H_t(z)\ne0$ for all $z\in R$ and $0\le t\le t_0$.
\end{theorem}

Theorem~\ref{thm:barrier} implies hypothesis (iii) of
Theorem~\ref{thm:criterion}: the curved barrier region of that
hypothesis satisfies, horizontally,
$X\le x\le X+\sqrt{1-y_0^2}\le X+1$; and vertically,
$y\ge\sqrt{y_0^2+2(t_0-t)}\ge y_0>\frac{1809}{10000}$ (indeed
$y_0^2-\bigl(\frac{1809}{10000}\bigr)^2=\frac{234599}{10^8}>0$) and
$y\le\sqrt{1-2t}\le1$; hence it is contained in $R\times[0,t_0]$.

\subsection{Uniform effective error on the box, including $t=0$}
\label{sec:barriererr}

\begin{proposition}[Uniform error bound on the barrier box]\label{prop:barriererr}
For all $x\in[X,X+1]$ and $0\le t\le t_0$, the window index satisfies
$N=\bigl\lfloor\sqrt{\frac{x}{4\pi}+\frac t{16}}\bigr\rfloor=690988$
and $N^2<x/(4\pi)$. Moreover, for all $z\in R$ and $0<t\le t_0$, the
error terms of Theorem~1.3, bounded with the constants of
Section~\ref{sec:weld}, satisfy
\begin{equation}\label{eq:barriererror}
e_A+e_B+e_{C,0}\ <\ 0.000356523011600040\ <\ 0.00125 .
\end{equation}
\end{proposition}

\begin{proof}
The proof is computer-assisted; the verification program
\texttt{prop62\_proof.c}, included with this manuscript as
supplementary material (directory \texttt{code/prop62}; derived
from Gomila's source as documented in its header and reviewed line
by line as part of this audit), encloses every quantity below in
$256$-bit ball arithmetic starting from the exact inputs $X$,
$t_0=\frac{129}{800}$, $y_{\min}=\frac{1809}{10000}$ and
$N=690988$; it reads no stored data, and every acceptance
condition is a strict ball inequality, which fails closed.

\emph{Window-index constancy.} The quantity
$q(x,t)=\frac{x}{4\pi}+\frac t{16}$ is strictly increasing in each
of $x$ and $t$, and $N=\lfloor\sqrt{q}\,\rfloor$ is nondecreasing
in $q$; so it suffices to evaluate $N$ at the two extreme corners
$(X,0)$ and $(X+1,t_0)$ of the box. At each corner the program
produces an enclosure of $\lfloor\sqrt q\rfloor$ that collapses to
an exact integer (a straddling enclosure would fail closed) and
checks that both integers equal $690988$. It further certifies
$N^2<X/(4\pi)\le x/(4\pi)$, and the equivalent inequality
$2\log N<L_{\min}:=\log\frac X{4\pi}$; write also
$L_{\max}:=\log\frac{X+1}{4\pi}$.

\emph{Error bound.} With $N$ constant, for $0<t\le t_0$ the error
terms are bounded by equation (23) of \cite{Polymath} and by
\eqref{eq:ec0}; the program encloses the following majorants, each
uniform on the box because every $t$- and $y$-dependent factor is
evaluated at its worst endpoint:
(i) by \eqref{eq:resbound}, using $(\,\cdot\,)_+\le1+8/X^2$ and
$y\ge y_{\min}$,
$\Rea s_*\ge a+\frac t4L_{\min}-\delta$ with
$a=\frac{1+y_{\min}}2$ and
$\delta=\frac{t_0(1+8/X^2)}{2X^2}$; since
$\log n\le\log N<\frac12L_{\min}$, this gives
$b_t(n)\,n^{-\Rea s_*}\le n^{-a_{\mathrm{eff}}}$ with
$a_{\mathrm{eff}}=a-\delta$, uniformly in $t\in[0,t_0]$ (the
$t$-dependent parts cancel against $\frac t4\log^2n$; certified:
$0<a_{\mathrm{eff}}<1$);
(ii) hence, by comparison with $\int_1^Nu^{-a_{\mathrm{eff}}}\,du$,
$\sum_{n=1}^{N}b_t(n)\,n^{-\Rea s_*}\le
1+\frac{N^{1-a_{\mathrm{eff}}}-1}{1-a_{\mathrm{eff}}}$;
(iii) since $n^2\le N^2<x/(4\pi)$, the bounds \eqref{eq:gammabound}
and \eqref{eq:kappabound} give
$|\gamma|n^y\le e^{0.02y}\bigl(\frac{4\pi n^2}x\bigr)^{y/2}\le
e^{0.02}$ and
$N^{|\kappa|}\le\exp\bigl(\frac{t_0\log N}{2(X-6)}\bigr)$, so the
prefactor of (23) is at most
$1+e^{0.02}\exp\bigl(\frac{t_0\log N}{2(X-6)}\bigr)$;
(iv) since $0<\log\frac x{4\pi n^2}\le L_{\max}$, the last factor
of (23) is at most
$\exp\bigl(\frac{(t_0^2/16)L_{\max}^2+0.626}{X-6.66}\bigr)-1$;
(v) in \eqref{eq:ec0}, the prefactor is maximized at
$(x,y)=(X,y_{\min})$, the factor $e^{-\frac t{16}\log^2\frac x{4\pi}}$
is at most $1$, $3^y+3^{-y}\le3+\frac13=\frac{10}3$ (increasing in
$y$, and $y\le1$), and the final quotient is bounded by maximizing
its numerator ($L\le L_{\max}$) and minimizing its denominator
($x\ge X$) separately.
The program multiplies and sums these enclosures --- yielding
$e_A+e_B\le3.64\times10^{-10}$ and
$e_{C,0}\le0.000356522648$ --- and certifies the two strict
inequalities of \eqref{eq:barriererror}, the sharp constant being
compared as the exact rational $356523011600040\times10^{-18}$.
The run completes in under a second; its output is included as
\texttt{code/prop62/runs/prop62\_report.txt}.
\end{proof}

The constancy of $N$ makes $f_t$ a single analytic
expression on $R$. The round number $0.00125=\frac1{800}$ is the
\emph{approximation allowance} used by the prism criterion below;
\eqref{eq:barriererror} says it is safely conservative.

Theorem~1.3 is stated in \cite{Polymath} for $t>0$; the following
lemma extends its conclusion to the closed endpoint $t=0$ on the
compact box.

\begin{lemma}[Extension to $t=0$]\label{lem:tzero}
For all $z\in R$ and all $0\le t\le t_0$ (including $t=0$),
$B_t(z)\neq0$ and
\[
\bigl|g_t(z)-f_t(z)\bigr|\ \le\ 0.00125,
\]
where $g_t=H_t/B_t$ and $f_t$ is the sum \eqref{eq:ft} with the
constant window index $N=690988$.
\end{lemma}

\begin{proof}
For $0<t\le t_0$ this is Theorem~1.3 together with
Proposition~\ref{prop:barriererr}. For the endpoint, fix $z\in R$
and argue by a limit $t_j\downarrow0$: (1)~by dominated convergence
in \eqref{eq:Ht} (majorant $e^{t_0u^2+u}|\Phi(u)|$, integrable by
the super-exponential decay of $\Phi$), $H_t(z)$ is jointly
continuous down to $t=0$ on $R$; (2)~$B_t^{\pm1}$ is continuous and
$B_t\neq0$ on $R\times[0,t_0]$, because the argument
$\frac{1+y-ix}{2}$ of $M_0$ stays away from the branch cut
(Section~\ref{sec:thm13}) and the deforming exponential never
vanishes; (3)~since the window index is constant, $f_t$ is a fixed
finite sum of functions continuous in $t$; (4)~the error majorant is
continuous in $t$ and bounded by $0.00125$ on $(0,t_0]$ by
\eqref{eq:barriererror}. Passing to the limit in
$|g_{t_j}-f_{t_j}|\le0.00125$ gives the claim at $t=0$.
\end{proof}

\subsection{Uniform derivative bounds on the box}\label{sec:derivbox}

The moving-mesh argument below needs uniform bounds for
$|\partial_zf_t|$ and $|\partial_tf_t|$ on $R$. The pointwise input
is the following result; recall that by
Proposition~\ref{prop:barriererr} the window index is constantly
$N=690988$ on the box, so that $f_t$ is the fixed finite sum
\eqref{eq:ft} with this $N$, jointly smooth in $(z,t)$, and the
lemma applies to it verbatim.

\begin{lemma}[{\cite[Lemma~8.4]{Polymath}}]\label{lem:polyderiv}
In the region \eqref{eq:region}, and away from the jump
discontinuities of $N$,
\begin{align*}
\Bigl|\frac{\partial f_t}{\partial z}\Bigr|
&\le\sum_{n=1}^{N}\frac{b_t(n)}{n^{\Rea s_*}}
\Bigl(\frac{\log n}2+\frac{t\log n}{4(x-6)}\Bigr)\\
&\quad+|\gamma|N^{|\kappa|}\sum_{n=1}^{N}
\frac{b_t(n)\,n^y}{n^{\Rea s_*}}
\Bigl(\frac{t\log n}{4(x-6)}
+\Bigl(\log\frac{|1+y+ix|}{4\pi}+\pi+\frac3x\Bigr)
\Bigl(\frac12+\frac t{4(x-6)}\Bigr)\Bigr),
\end{align*}
and
\begin{align*}
\Bigl|\frac{\partial f_t}{\partial t}\Bigr|
&\le\sum_{n=1}^{N}\frac{b_t(n)}{n^{\Rea s_*}}
\Bigl(\frac14\log n\log\frac x{4\pi n}+\frac\pi8\log n
+\frac{2\log n}{x-6}\Bigr)\\
&\quad+|\gamma|N^{|\kappa|}\sum_{n=1}^{N}
\frac{b_t(n)\,n^y}{n^{\Rea s_*}}
\Bigl(\frac14\log n\log\frac x{4\pi n}+\frac\pi8\log n
+\frac{2\log n}{x-6}
+\frac14\Bigl(\frac\pi2+\frac8{x-6}\Bigr)
\Bigl(\log\frac x{4\pi}+\frac8{x-6}\Bigr)\Bigr).
\end{align*}
\end{lemma}

To convert these pointwise bounds into uniform ones we introduce the
following frozen quantities. Write $y_{\min}=\frac{1809}{10000}$
(the bottom of $R$), $q(x)=\frac x{4\pi}$, and, as in
\eqref{eq:resbound},
\[
h(x,y)=1-3y+\frac{4y(1+y)}{x^2},
\qquad
c(x,y)=\frac14\log q(x)-\frac{h(x,y)_+}{2x^2},
\qquad
c_0=c(X,y_{\min}),
\]
so that \eqref{eq:resbound} reads
$\Rea s_*\ge\frac{1+y}2+t\,c(x,y)$. Define further
\begin{align*}
a(t) &= e^{0.02y_{\min}}\,q(X)^{-y_{\min}/2}\,
N^{\,t\,y_{\min}/(2(X-6))},
\\
P_t(u) &= u^{-\frac{1+y_{\min}}2+t(\frac14\log u-c_0)},
\\
\theta(t) &= \frac12+\frac t{4(X-6)},
\\
\Lambda_z(t) &= \Bigl(\log\frac{|2+i(X+1)|}{4\pi}+\pi+\frac3X\Bigr)
\theta(t),
\\
\Lambda_t &= \frac14\Bigl(\frac\pi2+\frac8{X-6}\Bigr)
\Bigl(\log\frac{X+1}{4\pi}+\frac8{X-6}\Bigr),
\\
C_2 &= \frac14\log\frac{X+1}{4\pi}+\frac\pi8+\frac2{X-6},
\end{align*}
and the two summand majorants
\begin{align*}
F_z(u;t)&=P_t(u)\Bigl[\theta(t)\log u
+a(t)\,u^{y_{\min}}\Bigl(\frac{t\log u}{4(X-6)}
+\Lambda_z(t)\Bigr)\Bigr],\\
F_t(u;t)&=P_t(u)\Bigl[\bigl(1+a(t)\,u^{y_{\min}}\bigr)
\log u\,\Bigl(C_2-\frac{\log u}4\Bigr)
+a(t)\,u^{y_{\min}}\Lambda_t\Bigr],
\end{align*}
and finally
\[
D_z(t)=\sum_{n=1}^{16}F_z(n;t)+\int_{16}^{N}F_z(u;t)\,du,
\qquad
D_t(t)=\sum_{n=1}^{16}F_t(n;t)+\int_{16}^{N}F_t(u;t)\,du .
\]

Two tight inequalities at the exact barrier parameters are needed
below. The first, $N^2<q(X)$, is part of
Proposition~\ref{prop:barriererr}. The second we record as a
standalone statement.

\begin{proposition}[Certified tail-exponent inequality]\label{prop:tailexp}
At $X=6000000185827$, $N=690988$, $y_{\min}=\frac{1809}{10000}$,
\[
\frac14\log\frac{q(X)}{N^2}\ >\ \frac{h(X,y_{\min})_+}{2X^2} .
\]
\end{proposition}

\begin{proof}
The proof is computer-assisted; the verification program
\texttt{prop65\_proof.c}, included with this manuscript as
supplementary material (directory \texttt{code/prop65}; written for
this review as a standalone extraction of the corresponding runtime
check of Gomila's barrier program, as documented in its header),
encloses both sides in $256$-bit ball arithmetic from the exact
rational parameters and certifies the displayed strict inequality
as a strict ball comparison, which fails closed. The certified
enclosures are
\[
\tfrac14\log\tfrac{q(X)}{N^2}
=2.2405813008\ldots\times10^{-7}\pm10^{-31},
\qquad
\frac{h(X,y_{\min})_+}{2X^2}
=6.3513884954\ldots\times10^{-27}\pm10^{-51};
\]
the inequality is tight only in the sense that its left side is
small --- it is positive precisely because $N^2<q(X)$, an
inequality of relative size $3\times10^{-7}$ --- while the right
side is smaller by twenty orders of magnitude. The run completes in
milliseconds; its output is included as
\texttt{code/prop65/runs/prop65\_report.txt}. The same inequality
is also re-checked at run time by the barrier program of
Proposition~\ref{prop:prismcert} (harmlessly, like the conditions
(C2), (C4) and (C7) of Proposition~\ref{prop:derivgates}).
\end{proof}

The construction consumes the following list of strict inequalities
at the exact barrier parameters.

\begin{proposition}[Derivative-bound conditions]\label{prop:derivgates}
At $X=6000000185827$, $N=690988$, $y_{\min}=\frac{1809}{10000}$,
$t_0=\frac{129}{800}$:
\begin{enumerate}
\item[\textup{(C1)}] $X^2>8$ and
$\frac1{4(X+1)}-\frac2{X^3}>0$;
\item[\textup{(C2)}] $0.02-\frac12\log q(X)
+\frac{t_0\log N}{2(X-6)}+\frac12\log N<0$;
\item[\textup{(C3)}] $2\log N<\log q(X)$;
\item[\textup{(C4)}] $\frac{12}{X^2}<3$;
\item[\textup{(C5)}] $c_0>\frac12\log N$;
\item[\textup{(C6)}] $\frac{1-y_{\min}}2>\frac1{\log16}$;
\item[\textup{(C7)}] $C_2>\frac14\log N$.
\end{enumerate}
\end{proposition}

\begin{proof}
(C1), (C4) and (C6) are immediate by inspection
($X\approx6\times10^{12}$;
$\frac1{\log16}=0.3607\ldots<0.40955=\frac{1-y_{\min}}2$). For
(C2), rewrite the left side as
$0.02+\frac12\log\frac{4\pi N}{X}
+\frac{t_0\log N}{2(X-6)}$: here
$\frac{4\pi N}{X}<\frac{8.8\times10^6}{6\times10^{12}}
<2\times10^{-6}<e^{-13}$, so the middle term is less than
$-\frac{13}2$, while the last term is less than $10^{-12}$; hence
the left side is less than $0.02-6.5+10^{-12}<0$. (C7) follows
from (C3): $C_2-\frac14\log N\ge\frac14\log q(X)-\frac14\log N
>\frac14\log N>0$.

The remaining two conditions are exactly the statements
established earlier: (C3) is the inequality $N^2<q(X)$ of
Proposition~\ref{prop:barriererr}, and (C5) is
Proposition~\ref{prop:tailexp}, since
$c_0-\frac12\log N
=\frac14\log\frac{q(X)}{N^2}-\frac{h(X,y_{\min})_+}{2X^2}$.
(The barrier program also re-checks (C2), (C4) and (C7) at run
time; by the above, those checks are mathematically redundant,
though harmless.)
\end{proof}

\begin{lemma}[Pointwise domination on the box]\label{lem:dompoint}
Assume (C1), (C2), (C4). Then for all $x\in[X,X+1]$,
$y\in[y_{\min},1]$, $0\le t\le t_0$ and $1\le n\le N$, the $n$-th
summand of the first (resp.\ second) bound of
Lemma~\ref{lem:polyderiv} is at most $F_z(n;t)$ (resp.\
$F_t(n;t)$). Consequently, for $0\le t\le t_0$,
\[
\sup_{z\in R}\Bigl|\frac{\partial f_t}{\partial z}(z)\Bigr|
\le\sum_{n=1}^{N}F_z(n;t),
\qquad
\sup_{z\in R}\Bigl|\frac{\partial f_t}{\partial t}(z)\Bigr|
\le\sum_{n=1}^{N}F_t(n;t).
\]
\end{lemma}

\begin{proof}
\emph{Step 1: $c$ is minimized at $(X,y_{\min})$.} In $y$:
$\partial_yh=-3+\frac{4(1+2y)}{x^2}\le-3+\frac{12}{X^2}<0$ by (C4),
so $h$, hence $h_+$, is nonincreasing in $y$ and $c$ is
nondecreasing in $y$. In $x$: where $h>0$,
\[
\Bigl|\partial_x\frac{h}{2x^2}\Bigr|
=\Bigl|-\frac{4y(1+y)}{x^5}-\frac{h}{x^3}\Bigr|
\le\frac{8}{x^5}+\frac{h_+}{x^3}\le\frac2{x^3},
\]
using $4y(1+y)\le8$, $8/x^2\le1$ and $h_+\le1$ (both from
$X^2>8$; for $y>0$, $h\le1$ reduces to $4(1+y)\le3x^2$, and at
$y=0$, $h=1$). Where $h<0$ the positive part is locally constant,
and $h_+/(2x^2)$ is locally Lipschitz across the kink with the same
bound. Hence for $x\in[X,X+1]$,
$\partial_xc\ge\frac1{4x}-\frac2{X^3}
\ge\frac1{4(X+1)}-\frac2{X^3}>0$ almost everywhere by (C1), so $c$
is increasing in $x$. Therefore $c(x,y)\ge c(X,y_{\min})=c_0$ on
the box.

\emph{Step 2: the $B$-side.} By \eqref{eq:resbound} and Step 1,
$\Rea s_*\ge\frac{1+y}2+t\,c(x,y)\ge\frac{1+y_{\min}}2+t\,c_0$, so
\[
b_t(n)\,n^{-\Rea s_*}
\le n^{\frac t4\log n-\frac{1+y_{\min}}2-t\,c_0}=P_t(n),
\]
while the $B$-side brackets obey
$\log n\bigl(\frac12+\frac t{4(x-6)}\bigr)\le\theta(t)\log n$ and
\[
\log n\Bigl(\frac14\log\frac x{4\pi n}+\frac\pi8+\frac2{x-6}\Bigr)
=\log n\Bigl(\frac14\log q(x)-\frac{\log n}4+\frac\pi8
+\frac2{x-6}\Bigr)
\le\log n\Bigl(C_2-\frac{\log n}4\Bigr),
\]
since $q$ is increasing and $\frac1{x-6}$ decreasing on $[X,X+1]$.

\emph{Step 3: the $A$-core.} By
\eqref{eq:gammabound}--\eqref{eq:kappabound},
$|\gamma|N^{|\kappa|}n^y\,b_t(n)\,n^{-\Rea s_*}\le
b_t(n)\,\mathcal A(x,y,t;n)$ with
\[
\mathcal A(x,y,t;n)
=e^{0.02y}q(x)^{-y/2}N^{\frac{ty}{2(x-6)}}\,
n^{\,y-\frac{1+y}2-t\,c(x,y)} .
\]
Using $\partial_yc\ge0$ (Step 1) and $\log n\ge0$,
\[
\partial_y\log\mathcal A
=0.02-\frac12\log q(x)+\frac{t\log N}{2(x-6)}
+\frac{\log n}2-t\,(\partial_yc)\log n
\le0.02-\frac12\log q(X)+\frac{t_0\log N}{2(X-6)}+\frac{\log N}2,
\]
which is negative by (C2); so $\mathcal A$ is maximized at
$y=y_{\min}$. At $y=y_{\min}$, freezing $x$ conservatively
($q(x)^{-y_{\min}/2}\le q(X)^{-y_{\min}/2}$,
$N^{\frac{ty_{\min}}{2(x-6)}}\le N^{\frac{ty_{\min}}{2(X-6)}}$,
$c(x,y_{\min})\ge c_0$) gives
$b_t(n)\,\mathcal A\le a(t)\,n^{y_{\min}}P_t(n)$.

\emph{Step 4: the $A$-side brackets.} On $[X,X+1]\times[y_{\min},1]$, we have
\begin{align*}
\frac{t\log n}{4(x-6)} &\le \frac{t\log n}{4(X-6)},
\\
\Bigl(\log\frac{|1+y+ix|}{4\pi}+\pi+\frac3x\Bigr)
\Bigl(\frac12+\frac t{4(x-6)}\Bigr)
&\le \Bigl(\log\frac{|2+i(X+1)|}{4\pi}+\pi+\frac3X\Bigr)
\theta(t)=\Lambda_z(t),
\end{align*}
since $|1+y+ix|\le|2+i(X+1)|$ (both coordinates increase) and
$\frac3x\le\frac3X$; and
\[
\frac14\Bigl(\frac\pi2+\frac8{x-6}\Bigr)
\Bigl(\log q(x)+\frac8{x-6}\Bigr)\le\Lambda_t,
\]
each factor being bounded by its value at the appropriate endpoint.

Combining Steps 2--4 termwise bounds the two summands of
Lemma~\ref{lem:polyderiv} by $F_z(n;t)$ and $F_t(n;t)$. For
$0<t\le t_0$ the displayed conclusions are then
Lemma~\ref{lem:polyderiv}; at $t=0$ they follow by continuity, as
both sides of each inequality are continuous in $t$ on the compact
box ($f_t$ being the fixed finite sum).
\end{proof}

\begin{lemma}[Head-and-tail evaluation]\label{lem:headtail}
Assume (C5), (C6), (C7). Then for each fixed $t\in[0,t_0]$ the
functions $u\mapsto F_z(u;t)$ and $u\mapsto F_t(u;t)$ are positive
and nonincreasing on $[16,N]$, and consequently
\[
\sum_{n=1}^{N}F_z(n;t)\ \le\ D_z(t),
\qquad
\sum_{n=1}^{N}F_t(n;t)\ \le\ D_t(t).
\]
\end{lemma}

\begin{proof}
Write $r=\log u\in[\log16,\log N]$. The logarithmic derivative of
$P_t$ with respect to $r$ is
$p(r)=-\frac{1+y_{\min}}2+t\bigl(\frac r2-c_0\bigr)
\le-\frac{1+y_{\min}}2$, since $t\ge0$ and
$\frac r2\le\frac12\log N\le c_0$ by (C5). Each of $F_z$, $F_t$ is
a nonnegative linear combination (with $u$-independent
coefficients $\theta(t)$, $\frac t{4(X-6)}$, $a(t)$,
$\Lambda_z(t)$, $\Lambda_t$, all $\ge0$) of the six component
functions
\[
P_t(u)\,r,\quad P_t(u)\,u^{y_{\min}}r,\quad
P_t(u)\,u^{y_{\min}},\quad
P_t(u)\,B(r),\quad P_t(u)\,u^{y_{\min}}B(r),
\]
with $B(r)=r\,(C_2-\frac r4)$, which is positive on
$[\log16,\log N]$ by (C7). Their logarithmic $r$-derivatives are,
respectively, at most
\[
p+\frac1r,\quad p+y_{\min}+\frac1r,\quad p+y_{\min},\quad
p+\frac{B'}B,\quad p+y_{\min}+\frac{B'}B,
\]
and $\frac{B'}B=\frac1r-\frac1{4(C_2-r/4)}\le\frac1r$. Hence every
component has logarithmic derivative at most
\[
p+y_{\min}+\frac1{\log16}
\le-\frac{1+y_{\min}}2+y_{\min}+\frac1{\log16}
=-\frac{1-y_{\min}}2+\frac1{\log16}\ <\ 0
\]
by (C6), so each component, and therefore each of $F_z,F_t$, is
positive and decreasing on $[16,N]$. Finally, for decreasing
positive $F$,
$\sum_{n=17}^{N}F(n)\le\int_{16}^{N}F(u)\,du$ (compare each term
with the integral over $[n-1,n]$), which gives the displayed
bounds.
\end{proof}

\begin{lemma}[Uniform derivative bounds]\label{lem:derivbox}
Assume (C1)--(C7). Then for every $t\in[0,t_0]$,
\[
\sup_{z\in R}\Bigl|\frac{\partial f_t}{\partial z}(z)\Bigr|
\ \le\ D_z(t),
\qquad
\sup_{z\in R}\Bigl|\frac{\partial f_t}{\partial t}(z)\Bigr|
\ \le\ D_t(t),
\]
and hence, for every closed interval
$[t_i,t_{i+1}]\subseteq[0,t_0]$, the relation
$$\sup_{z\in R,\ \tau\in[t_i,t_{i+1}]}|\partial_\tau f_\tau(z)|
\le\sup_{\tau\in[t_i,t_{i+1}]}D_t(\tau)$$
holds.
\end{lemma}

\begin{proof}
Combine Lemmas~\ref{lem:dompoint} and~\ref{lem:headtail}.
\end{proof}

\begin{remark}\label{rem:quadrature}
The barrier program evaluates $D_z(t)$ and $D_t(t)$ with $t$ an
interval enclosure --- for the prism bound, an enclosure of the
whole interval $[t_i,t_{i+1}]$, which by inclusion isotonicity
encloses $\sup_\tau D_t(\tau)$ --- and computes the integrals
$\int_{16}^NF(u)\,du$ by Arb's rigorous Gauss--Legendre
integrator~\cite{JohanssonIntegration}. That algorithm requires the
integrand callback to supply enclosures that are holomorphic on the
complex balls it queries; the integrands above extend
holomorphically in $u$ away from the branch cut of the logarithm,
and the program uses the analytic logarithm and power operations,
which return non-finite enclosures (failing the run) if a queried
ball touches the cut. Two interfaces of the original Polymath-era
barrier code are thereby repaired: the discrete sums are no longer
bounded by the unjustified expression $F(1)+\int_1^NF$ (the summand
is not monotone near $n=1$; Lemma~\ref{lem:headtail} certifies
monotonicity only from $n=16$ on), and the callbacks no longer
project the integration variable to its real part, which would void
the integrator's error analysis.
\end{remark}

\subsection{The prism criterion}\label{sec:prism}

The time interval $[0,t_0]$ will be covered by finitely many closed
``prisms'' $R\times[t_i,t_{i+1}]$. The following lemma reduces
nonvanishing of $H$ on a prism to finitely many verifiable
conditions at its initial time $t_i$.

\begin{lemma}[Prism criterion]\label{lem:prism}
Let $0\le t_i<t_{i+1}\le t_0$, and let $z_1,\dots,z_m$
(cyclically ordered, $z_{m+1}=z_1$) be a mesh on $\partial R$ with
consecutive spacing at most $h$. Let $D_z\ge\sup_{z\in R}|\partial_zf_{t_i}(z)|$ and
$D_t\ge\sup_{z\in R,\ \tau\in[t_i,t_{i+1}]}
|\partial_\tau f_\tau(z)|$ (as supplied by
Lemma~\ref{lem:derivbox}), and suppose:
\begin{enumerate}
\item[\textup{(a)}] $\bigl|f_{t_i}(z_j)\bigr|\ge M_i$ for all $j$,
where
\begin{equation}\label{eq:gate}
M_i\ >\ \frac{D_{z}h}{2}
+D_{t}\,(t_{i+1}-t_i)+0.00125\,;
\end{equation}
\item[\textup{(b)}] each argument increment
$\arg\bigl(f_{t_i}(z_{j+1})/f_{t_i}(z_j)\bigr)$ lies strictly in
$(-\pi,\pi)$, and the total winding of the polygon
$\bigl(f_{t_i}(z_j)\bigr)_j$ about $0$ is~$0$.
\end{enumerate}
Then $H_\tau(z)\neq0$ for all $z\in R$ and
$\tau\in[t_i,t_{i+1}]$.
\end{lemma}

\begin{proof}
On each half of a mesh subedge, both the true image curve
$F(s)=f_{t_i}(\text{subedge}(s))$ and its endpoint chord lie in the
closed disk of radius $D_{z}h/2$ about the nearer endpoint value
(this is the factor $\frac12$ in \eqref{eq:gate}); each disk is
convex, so the straight-line homotopy from the true boundary curve
to the mesh polygon stays inside these disks. Moving to time
$\tau\in[t_i,t_{i+1}]$ displaces values by at most
$D_{t}(\tau-t_i)$, and passing from $f_\tau$ to
$g_\tau=H_\tau/B_\tau$ by at most $0.00125$
(Lemma~\ref{lem:tzero}). Condition \eqref{eq:gate} says every
resulting disk, centered at a mesh value of modulus $\ge M_i$,
excludes the origin; hence the entire three-stage homotopy (polygon
$\to$ true $f_{t_i}$-boundary image $\to$ $f_\tau(\partial R)$ $\to$
$g_\tau(\partial R)$) avoids zero. In particular each chord avoids
the origin, so its argument variation is the principal argument of
the ratio of its endpoint values and the winding number of the
polygon about the origin is a well-defined integer; by (b) it is
$0$, and the zero-avoiding homotopy gives winding number $0$ for
$g_\tau(\partial R)$ as well, for every $\tau\in[t_i,t_{i+1}]$.
Since $B_\tau\neq0$ and $B_\tau$ is holomorphic on a neighborhood
of $R$ (Lemma~\ref{lem:tzero} and Section~\ref{sec:thm13}),
$g_\tau$ is holomorphic there and has the same zeros as $H_\tau$,
and $g_\tau$ is zero-free on $\partial R$; the argument principle
then gives $H_\tau\neq0$ on $R$.
\end{proof}

\subsection{The certified prism decomposition}\label{sec:prismcert}

The computational input of the barrier argument is the following.

\begin{proposition}[Certified prism decomposition]\label{prop:prismcert}
There exist times $0=\tau_0<\tau_1<\dots<\tau_{883}$ with
$\tau_{883}\ge t_0$ such that for each $0\le i<883$, the
hypotheses (a)--(b) of Lemma~\ref{lem:prism} hold for the prism
$R\times[\tau_i,\min(\tau_{i+1},t_0)]$, with certified margin in
\eqref{eq:gate} at least $0.5198498946$ in every prism.
\end{proposition}

\begin{proof}
The proof is computer-assisted; the verification program
\texttt{prop612\_proof.c}, included with this manuscript as
supplementary material (directory \texttt{code/prop612}; assembled
for this review from three components of Gomila's barrier pipeline,
each descended from the Polymath codebase, as documented in its
header and reviewed line by line as part of this audit), is fully
self-contained: it reads no stored data, and every decisive
comparison is a strict Arb ball comparison, which fails closed. It
runs in three phases.

\emph{(1) Domain hypotheses and truncation.} The program first
certifies, at $192$-bit precision: the uniform expansion-parameter
bound $|w|\le0.6480<0.66$ on the whole slab (via the shift bounds
$|\Rea(\log n_0-\alpha)|<0.694$ and $|\Ima(\log n_0-\alpha)|<0.7855$
for both $\alpha$-arguments used below, where $n_0=N/2$); the bound
$0.0882<0.089$ for the modulus of the $\gamma$-type factor, both
for its exact expression (on a $256\times128$ subdivision of the
box) and for the implemented majorant; and, by a rigorous direct
sum over $n\le N$ with geometric tail factors, the complete
two-variable $62$-term Taylor truncation bound
$1.9543\times10^{-22}<10^{-20}$ --- the allowance that is attached
as an error ball to every boundary mesh value in phase (3).

\emph{(2) Coefficient generation.} The $62\times62$ complex matrix
\[
\mathrm{finalmat}[e][k]
=\sum_{n=1}^{N}n^{\frac{i(X+\frac12)-1}{2}}
\,\frac{\bigl(\log\frac n{n_0}\bigr)^e}{e!}
\,\frac{\bigl(\frac14\log^2\frac n{n_0}\bigr)^k}{k!}
\]
is computed in memory by direct summation, each entry a genuine
ball enclosure. The reconstruction identity --- verified
symbolically in the course of this review --- shows that the
polynomial evaluation of phase (3) applied to this matrix
reproduces, exactly up to the certified truncation of phase (1),
the two sums of \eqref{eq:ft} at the evaluation point, with the
heat factor $b_t(n)$ emerging from the algebra of the two Taylor
directions. (In Gomila's pipeline this matrix was read from a
stored $20$-digit file and independently regenerated for
comparison; computing it in-program from the defining series
removes the stored file entirely and settles the provenance of the
coefficients.)

\emph{(3) The adaptive prism loop.} Before the loop the program
re-certifies the conditions (C1)--(C7) of
Proposition~\ref{prop:derivgates} at run time. Then, starting at
$\tau_0=0$: at each prism start $\tau_i$ it evaluates the
derivative bounds $D_z(\tau_i)$ and (for the candidate prism, with
the time parameter interval-valued over the entire closed prism)
$D_t$, by the head-plus-integral quadratures of
Section~\ref{sec:derivbox}, rounding both up to exact integers;
meshes $\partial R$ with $4\lceil D_z\rceil-4$ points (spacing at
most $h=1/(\lceil D_z\rceil-1)$); requires every mesh enclosure of
$f_{\tau_i}$ to exclude zero, every argument increment to lie
strictly in $(-\pi,\pi)$, and the per-rectangle winding enclosure
to lie strictly in $(-\frac14,\frac14)$, certifying the integer
winding number $0$; chooses $\tau_{i+1}$ adaptively at half the
certified slack of \eqref{eq:gate}, as an exact dyadic number; and
finally re-checks the complete gate
$\min_j|f_{\tau_i}(z_j)|>\frac{D_zh}2+D_t(\tau_{i+1}-\tau_i)
+0.00125$ as a strict ball inequality. The loop ends only when the
prisms cover $[0,t_0]$ exactly (adjacent seams are identical exact
dyadic numbers, the first time is exactly $0$, the final endpoint
encloses $\frac{129}{800}$), and the aggregate winding enclosure
must lie strictly in $(-\frac14,\frac14)$. (The final prism may
extend slightly beyond $t_0$; restricting it to its intersection
with $[0,t_0]$ only shrinks the displacement terms in
\eqref{eq:gate}, so the certificate for the full final prism covers
the truncated one, as required by the statement.)

The fresh run certifies $883$ prisms --- the same count as
Gomila's sealed transcripts --- with aggregate winding enclosure
$[\pm1.18\times10^{-14}]$ and worst certified margin
$0.51984989461387254365\pm2\times10^{-20}$, certified against the
exact rational threshold $5198498946\times10^{-10}$ (a check added
in this review; Gomila's original certified only positivity of each
margin). The run completes in under eight minutes; its output is
included as \texttt{code/prop612/runs/prop612\_report.txt}.
\end{proof}

\begin{proof}[Proof of Theorem~\ref{thm:barrier}]
By Proposition~\ref{prop:prismcert} and Lemma~\ref{lem:prism},
$H_\tau(z)\neq0$ for $z\in R$ and
$\tau\in[\tau_i,\min(\tau_{i+1},t_0)]$, for each $i$. The union of
these intervals is $[0,t_0]$.
\end{proof}

% ===============================================================
\section{Assembly: proof of Theorem \ref{thm:main}}
% ===============================================================

\begin{proof}[Proof of Theorem \ref{thm:main} (candidate argument)]
The parameters \eqref{eq:params} satisfy $0<t_0<\frac12$ and
$0<y_0^2<1-2t_0$. Hypothesis (i) of Theorem~\ref{thm:criterion}
holds by Section~\ref{sec:hypi}, since $X/2<T_{\rm PT}$. Hypothesis
(ii) holds by combining Proposition~\ref{prop:finitelane}
(nonvanishing of $H_{t_0}$ for $x_*\le x<x_{3840001}$,
$y_0\le y\le\sqrt{1-2t_0}$) with Theorem~\ref{thm:tail} applied at
$t=t_0\in I_t$ (nonvanishing for $x\ge x_{3840000}=X_{N_*}(t_0)$):
the two ranges overlap on the full window $W_{3840000}$, so together
they cover all $x\ge x_*$. Hypothesis (iii) holds by
Theorem~\ref{thm:barrier}, since the curved barrier region is
contained in $R\times[0,t_0]$
(Section~\ref{sec:barrier}). Theorem~\ref{thm:criterion} then gives
\[
\LB\ \le\ t_0+\frac{y_0^2}{2}\ =\ \frac{893927}{5000000}
\ =\ 0.1787854 . \qedhere
\]
\end{proof}

\section{Summary of the verification}
\label{sec:gaps}
% ===============================================================

This section records the status of the verification.

\emph{Published inputs.} The statements cited from \cite{Polymath}
and \cite{PlattTrudgian} have been checked verbatim against the
arXiv source files and are applied on domains where their
hypotheses hold, the one caveat being Remark~\ref{rem:1044} (the
constant $10.44$ of \cite[eq.~(24)]{Polymath} is not used; the
conservative $10.50$ corollary \eqref{eq:ec0} is derived and used
instead).

\emph{Analytic arguments.} Every analytic lemma carries a complete
proof verified in the course of the review: the proofs of
Lemmas~\ref{lem:frozen}, \ref{lem:contraction}, \ref{lem:tzero}
and~\ref{lem:prism} were verified in detail, and the formerly
outlined proofs of Lemma~\ref{lem:reduction},
Lemma~\ref{lem:derivbox} and Proposition~\ref{prop:transfer} were
completed in full, through
Lemmas~\ref{lem:capdecay}--\ref{lem:ydecay} and~\ref{lem:errmaj},
Lemmas~\ref{lem:dompoint} and~\ref{lem:headtail} with
Proposition~\ref{prop:derivgates}, and
Lemmas~\ref{lem:sigmaslope}--\ref{lem:cells} with
Proposition~\ref{prop:dinicells}, respectively. No analytic items
remain.

Two of these reductions have additionally been machine-checked. The
extension of a lower bound on $(0,t_0]$ to the closed endpoint
$t=0$ used in Lemma~\ref{lem:tzero}, and the reduction of the
zeros of $H_t$ to those of $H_t/B_t$ on a region where
$B_t\neq0$, are formalised in Lean~4 with \texttt{mathlib} in the
repository's \texttt{lean/} directory. The formalisation is
\texttt{sorry}-free and depends only on the three standard
\texttt{mathlib} axioms
(\texttt{propext}, \texttt{Classical.choice}, \texttt{Quot.sound});
in particular $B_t\neq0$ is obtained there structurally, from the
closed form of $M_0$, rather than from a numerical lower bound on
$\lVert M_0\rVert$. What is formalised is the reduction logic
together with the nonvanishing; the identification of the
formalised closed form with the implementation remains an ordinary
review item, as does the joint continuity of the true $H_t$ up to
$t=0$, whose Lean statement is in place but whose concrete
instantiation is not yet discharged.

\emph{Computer-assisted components.} Each of the 
propositions whose proofs are computer-assisted --- \ref{prop:sweep} (window sweep),
\ref{prop:dinicells} (Dini cells), \ref{prop:finiteerr}
(finite-region error bound), \ref{prop:tailcert} (tail cutoff
inequalities), \ref{prop:barriererr} (barrier-box error bound),
\ref{prop:tailexp} (tail-exponent inequality) and
\ref{prop:prismcert} (prism decomposition) --- has been
independently reproduced by a standalone program included as
supplementary material (directory \texttt{dan-reworking/code/}; the promoted copies are sealed under \texttt{independent/}), derived from
Gomila's sources as documented in each program's header, reviewed
line by line against the corresponding statement in this paper,
compiled against an independently built FLINT/Arb toolchain,
reading no stored data, and with every displayed constant of the
statements certified by an explicit strict comparison against an
exact rational. The proofs given with the propositions record the
protocols and the fresh-run results. In particular the stored
certificate files, transcripts and coefficient data of Gomila's
repository are nowhere consumed: every number entering
Theorem~\ref{thm:main} is either proved analytically in this paper
or recomputed from exact rational inputs by the supplementary
programs.

\subsection*{Overall assessment of proof architecture}
The theorem inputs are settled: the criterion, the effective
approximation, and the verified-height input are exactly as the
candidate uses them; every analytic argument carries a complete
verified proof; and every certified computation has been
independently reproduced as described above. The logical
architecture is clean: three independent suppliers feed the three
hypotheses of a single published criterion, the finite and tail
ranges overlap on a full window, the barrier rectangle strictly
contains the curved barrier, and every numerical acceptance
condition is a strict directed interval inequality that fails
closed. The thinnest quantitative margins are the finite-range
post-error margin ($5.6\times10^{-7}$), the height-transfer ratio
($1-\text{ratio}\approx1.4\times10^{-6}$), and the tail contraction
($1-D\approx2.8\times10^{-4}$); each was reproduced exactly, digit
for digit, on the independent toolchain. The remaining trust base
is the standard one for computer-assisted proofs of this kind: the
correctness of the FLINT/Arb library, of the compiler, and of the
line-by-line correspondence between the supplementary programs and
the statements of this paper, which the review has checked and
which any reader can re-audit from the included sources.

One caveat on re-execution, as distinct from re-reading. Six of the
seven programs can be re-run directly from what is supplied:
Propositions~\ref{prop:dinicells} and~\ref{prop:tailexp} take no
arguments, Proposition~\ref{prop:tailcert} takes an optional
precision defaulting to $320$ bits, and the remaining reports
record the parameters used. The exception is
Proposition~\ref{prop:sweep}, whose program takes twelve positional
arguments of which the stored sweep outputs record only the range,
the mollifier type, the $t$-box and the weight; the remaining
arguments are not recorded, so those nine runs cannot presently be
regenerated from the supplied material alone. The sources are
complete and auditable by reading; the invocations are not.

% ===============================================================
\begin{thebibliography}{9}

\bibitem{deBruijn}
N.~G. de Bruijn,
\emph{The roots of trigonometric integrals},
Duke Math.\ J.\ \textbf{17} (1950), 197--226.

\bibitem{CsordasNorfolkVarga}
G.~Csordas, T.~S. Norfolk and R.~S. Varga,
\emph{A lower bound for the de Bruijn--Newman constant $\Lambda$},
Numer.\ Math.\ \textbf{52} (1988), 483--497.

\bibitem{KKL}
H.~Ki, Y.-O. Kim and J.~Lee,
\emph{On the de Bruijn--Newman constant},
Adv.\ Math.\ \textbf{222} (2009), 281--306.

\bibitem{JohanssonIntegration}
F.~Johansson,
\emph{Numerical integration in arbitrary-precision ball
arithmetic},
Mathematical Software -- ICMS 2018, Lecture Notes in Comput.\ Sci.\
\textbf{10931}, Springer (2018), 255--263.
arXiv:1802.07942.

\bibitem{Newman}
C.~M. Newman,
\emph{Fourier transforms with only real zeros},
Proc.\ Amer.\ Math.\ Soc.\ \textbf{61} (1976), 245--251.

\bibitem{PlattTrudgian}
D.~Platt and T.~Trudgian,
\emph{The Riemann hypothesis is true up to $3\cdot10^{12}$},
Bull.\ Lond.\ Math.\ Soc.\ \textbf{53} (2021), 792--797.
arXiv:2004.09765.

\bibitem{Polymath}
D.~H.~J. Polymath,
\emph{Effective approximation of heat flow evolution of the Riemann
$\xi$ function, and a new upper bound for the de Bruijn--Newman
constant},
Res.\ Math.\ Sci.\ \textbf{6} (2019), no.\ 31.
arXiv:1904.12438v2.

\bibitem{RodgersTao}
B.~Rodgers and T.~Tao,
\emph{The de Bruijn--Newman constant is non-negative},
Forum Math.\ Pi \textbf{8} (2020), e6.

\end{thebibliography}

\end{document}
